% Modernised reading — fonds Alexandre Grothendieck, shelfmark 81
% (pages 1-54, the whole folder).
%
% This is an INTERPRETATION, not a transcription. It re-reads the pages in
% today's notation and vocabulary, states the hypotheses the manuscript leaves
% implicit, and drops the critical apparatus so that the mathematics reads
% straight through. Everything the brackets and underlines carried moves into a
% footnote. It is held to being mathematically correct as it stands; where the
% manuscript is loose or mistaken, the true statement is what appears here and
% the footnote says what the page has. In any dispute of substance the
% transcription governs: whoever cites Grothendieck must cite batch-01.fr.tex
% (pp. 1-20), batch-02.fr.tex (pp. 21-40) or batch-03.fr.tex (pp. 41-54).
%
% Pass: Opus 5.5 (claude-opus-5-5), 2026-09-23 — first pass, unchecked against the pages by a human.
% Revised 2026-09-23 (Opus 5.5): (1) no order in time the leaves do not give —
% the résumé no longer calls p. 2 the « point de départ » of the dessins
% d'enfants (it now says these are the objects the Esquisse d'un programme
% (1984) calls so, and that the pages cite neither Belyi nor the Esquisse);
% likewise « devenue célèbre », « postérieurs à la page » (names footnote,
% p. 2), « définira » (p. 26), « postérieur ou contemporain » (pocsets
% footnote), « connaissait / ait connu » (Buneman, pocsets, Kreweras
% footnotes) replaced by what the pages cite; ordinal « première / deuxième /
% troisième / dernière rédaction » replaced by the leaves' own labels
% (Version I, pp. 25-29, Version II); « nouveau départ », « reprise » (pp.
% 25, 51) and past tenses about Version I made neutral; the résumé's « avant
% d'être défini » (margin of p. 26 vs p. 36) and « quand une rédaction ne
% tient plus » removed. (2) First dictionary (p. 2): the bijection is now
% stated for the right classes — special functions (ramification <= 2 over
% 1) <-> marked maps with free half-edges; index exactly 2 over 1 <-> marked
% maps as defined (no free half-edge); edges and half-edges are closures of
% components of f^{-1}(]0,1]) (not of [0,1], whose preimage is the whole
% graph); the footnote says what the page has and gives T, T^2 as the
% failure of the loose version. Spine list item 1 and the résumé's converse
% adjusted to match.
% The folder was read whole, its three batches together; this is one document
% for the shelfmark. It was made on Opus 5.5 and has not been measured against
% the reference reading of folder 115 (made on Opus 5). The literature named
% in the footnotes is cited from memory and was not checked against the
% sources.
%
% Scope. Pages 1, 13 and 40 are covers whose words head sections here
% (« Cartes sphériques en général », « Décompositions n-aires (et découpages)
% d'espaces topologiques », « relations d'équi. non parallèles »). Pages 6-7
% are in another hand and unrelated; page 8 is blank; page 12 carries nothing
% for the mathematics; page 30 is blank. None of them is described. The
% inventory title (« Immersions du disque et de la circonférence ») is on no
% transcribed leaf and describes none; footnoted in the spine section.
%
% Spine. How pieces cut out by non-crossing cuts fit together, in three
% settings, with a map-theoretic frame at both ends:
%   pp. 1-5     spherical maps with a marking, free half-edges allowed <->
%               rational functions P/Q normalised at 0, 1, oo, unramified off
%               {0,1,oo}, ramification <= 2 over 1; maps without free
%               half-edges <-> exactly 2 (clean Belyi maps; dessins).
%               Coefficients algebraic, Aut(C) acts through Gal(Qbar/Q).
%               Pp. 3, 5 cancelled scratch.
%   pp. 9-11    drawings: discs cut by chords, with small graphs; p. 10 power
%               sets and Hom(X,Y) in P(X x Y).
%   pp. 13-24   Version I: closed cuts E_i of a connected space, each with two
%               complementary components; two cuts give five pieces; n cuts
%               give n+1 pieces, each X minus disjoint closed hemispheres,
%               uniquely; the pieces and cuts form a tree (stated p. 23).
%   pp. 25-29   « disposition binaire »; from a tree to a family of pairwise
%               (strictly) non-parallel binary decompositions, conditions
%               (*), (**), (***).
%   pp. 31-39   Version II, §§ 1.1-1.3: closed hemispheres, equator, the
%               non-parallel pair and diagram (14), strictness (21), the
%               ordered set C with involution sigma (25)-(27), the « quartette
%               ordonné », a Proposition (p. 39) whose proof breaks off.
%   pp. 40-52   polygons (bi-circulated finite sets), segments, non-parallel
%               parts = non-crossing, « partitions non parallèles » = the
%               non-crossing partitions of Kreweras (1972), nested families
%               of segments and cores, anglets, the same for a quasi-segment.
%   pp. 53-54   gluing two maps at an edge: faces and genus.
%
% Notation fixed for the document.
%  - Version I (pp. 14-24): Sigma_i^+-, the two OPEN components of X \ E_i;
%    hat-Sigma_i^+- = Sigma_i^+- u E_i, CLOSED. Version II (pp. 31-39):
%    Sigma^eps the CLOSED hemispheres, Sigma^{0 eps} = Sigma^eps \ E open.
%    So hat-Sigma (I) = Sigma (II) and Sigma (I) = Sigma^0 (II); stated in the
%    spine section, kept as each version writes it.
%  - « morceau » for a class of the partition Pi_I^0 (the page: « classe »).
%  - Gamma_I for the graph of p. 23 (the page: (S, I, R)); R is kept for the
%    equivalence relation of pp. 47-50 only. The residual set of p. 46 (the
%    page's R) is written Z.
%  - Polygon (S, {u, u^{-1}}), N = card S >= 3; « segment » = arc of
%    consecutive vertices, 0 < card < N; d T its two ends.
%  - « non parallèle » is the page's word throughout and is kept; it means
%    « non-crossing / compatible » (two disjoint chords in the usual sense
%    parallel are non parallèles in his). Said once, in the spine section.
%  Collisions not corrected: S is the set of pieces (p. 23), of tree
%  vertices (p. 27) and the polygon (p. 41); I indexes the cuts (pp. 14-24)
%  and is the quasi-segment (p. 51); A indexes a family (p. 37), is the edge
%  set of a tree (p. 27) and a part of the polygon (pp. 42 ff.).
%
% Substantive departures, each footnoted where it occurs:
%  p. 2    the first « ≃ » (maps <-> f étale off 0,1,oo) needs the
%          ramification <= 2 over 1 that the next line adds; without it, f
%          classifies hypermaps (bipartite maps); with it, maps with free
%          half-edges; maps without them <-> index exactly 2 over 1 (clean
%          Belyi maps), stated as two bijections. The page's question « P'=0
%          sans P=Q ? » answered yes; in the family aT^2+bT only T^2 and
%          2T-T^2 are special (ours).
%  p. 4    algebraicity of the coefficients: argument supplied (finitely
%          many special f of each degree).
%  p. 10   « Hom(E, Hom(X,Z)) = P(Hom(X,Z)) »: variance, P(Y) = Hom(Y, E).
%  p. 14   X connected and I finite supplied; the components are open without
%          local connectedness (two of them); pieces need not be connected
%          (example on the circle, ours).
%  pp. 16, 17  « X \ Sigma_i » read X \ E_i (sic twice).
%  p. 16   the equality for closure(Sigma_12) proved (ours).
%  p. 22   the Lemma needs (b), not disjointness alone; proof supplied.
%  p. 23   b) read as « no multiple edges »; the tree Lemma proved (ours).
%  p. 24   « K \ {i} = {alpha | E_alpha in Sigma_alpha^-} » is wrong; with
%          Sigma_i^- it fails for three parallel chords; correct form ours.
%  p. 27   A', A'' must contain a_0; Gamma finite.
%  p. 32   (11), (11 bis) are implied by (10), not equivalent to it; the
%          equivalent forms use the open Sigma^{0-} (counterexample ours).
%  p. 36   the Cor.'s proof: Sigma_1^- in Sigma_1^+ gives Sigma_1^+ = X.
%  p. 38   transitivity of (26), asked in the margin, checked; (27) proved.
%          The quartette's axiom a) does not imply that sigma reverses the
%          order (counterexample ours): the NB is an axiom.
%  p. 39   the Proposition is FALSE as stated: three chords around a triangle
%          (the tripod of p. 11) give C_{>= x} not totally ordered; true:
%          every interval [x, z] is a chain, and C_{>= x} is a chain for all
%          x exactly in the linear case. « (28 a) » identified with a) of
%          p. 38.
%  p. 43   the smallest pair (Abar, Bbar) supplied.
%  p. 46   « partition par-dessus » read as the non-parallel (non-crossing)
%          partition of p. 44.
%  p. 48   the Prop's proof, as legible, needs repair (A_i may leave A);
%          proof ours.
%  p. 49   « R serait un segment contenant deux anglets » read as « S_1
%          serait un segment saturé, contenant donc un anglet »; R_1 = S_1.
%  p. 53   the operation is not described on the page; reconstructed as the
%          exchange of the ends of two edges; mon/mon needs a connected
%          result (not two bridges).
%
% Links the pages do not write (ours, and said so in the body): Version I
% goes from cuts to a tree, pp. 27-29 from a tree to decompositions; the
% counterexample to p. 39 is the tripod drawn on p. 11; the three parallel
% chords of p. 47 and the four of p. 9; the nested families of p. 45 and the
% chord diagrams of p. 44; the bi-circulation of p. 41 is folder 89's third
% description of a polygon (89, pp. 19-20); p. 53 is the combinatorial side
% of folder 67's surgery and of folder 88's maps.
%
% Announced and never established: the end of p. 4 (« le diviseur de f
% est »); the question of the drawings pp. 9-11 (no text); the equality of
% p. 16 (stated); the Lemma of p. 22 (« Cela provient »); the tree Lemma of
% p. 23 (proof begun p. 24, stops); the « description complète » of p. 28;
% p. 29's « à examiner de plus près »; the transitivity asked on p. 38; the
% converse announced by « Inversement » (p. 38); the Proposition of p. 39;
% « il existe un plus petit… » (p. 43); the Intuition of p. 44; « Je veux
% construire A » (p. 47, refuted by its own NB); the tower of p. 50; the NB
% of p. 52 and the study of Sigma_max; the derivation of p. 53's counts.
\documentclass[11pt,a4paper]{article}
% Le préambule commun à toutes les transcriptions du fonds Grothendieck.
%
% Il tient en une idée : **l'incertitude se compose, elle ne se résout pas.**
% Une transcription qui lisse un mot douteux a détruit la seule chose pour
% laquelle on la faisait — un lecteur doit pouvoir savoir, à chaque endroit,
% ce qui était sur le feuillet et ce qui a été ajouté. Les six macros qui
% suivent sont donc l'appareil critique, et non de la décoration.
%
% Les mêmes commandes sont reconnues par `scripts/render.mjs`, qui produit la
% vue de lecture du site. Une macro ajoutée ici doit l'être aussi là-bas,
% faute de quoi le rendu échouera bruyamment — ce qui est voulu : un rendu
% silencieusement faux est pire qu'un rendu absent.

\ProvidesPackage{grothendieck}[2026/08/08 Transcriptions du fonds Grothendieck]

\usepackage{amsmath,amssymb,amsthm}
\usepackage{mathrsfs}
\usepackage{stmaryrd}
% Les diagrammes commutatifs sont partout dans ce fonds : c'est la langue
% même de la géométrie algébrique de Grothendieck, et une transcription qui
% les rendrait en prose ne transcrirait rien.
\usepackage{tikz-cd}
% Français par défaut — c'est la langue des feuillets — mais l'anglais est
% chargé aussi : la traduction et le résumé sont des documents anglais, et la
% typographie française (espaces avant les deux-points) y serait une faute.
% Ces documents basculent par \selectlanguage{english} en tête de corps.
\usepackage[english,french]{babel}
\usepackage{fontspec}
\usepackage{geometry}
\geometry{a4paper,margin=25mm,marginparwidth=28mm}
\usepackage{marginnote}
\usepackage{xcolor}
% ulem, not soul. `soul`'s \st and \ul are built on a character-by-character
% scan that XeLaTeX's font machinery breaks: the document fails with an
% undefined control sequence rather than degrading. ulem does the same two jobs
% and is Unicode-engine safe. `normalem` stops it hijacking \emph, which the
% transcriptions use for Grothendieck's own emphasis.
\usepackage[normalem]{ulem}
\usepackage{hyperref}
\hypersetup{colorlinks=true,linkcolor=black,urlcolor=black,pdfborder={0 0 0}}

% --- Métadonnées du lot -----------------------------------------------------
% Reprises telles quelles par le script de rendu, qui les affiche en tête de
% la vue de lecture. Elles fixent ce que le fichier prétend couvrir : sans
% elles, une transcription flotte sans point d'attache dans un fonds de seize
% mille feuillets.

\newcommand{\folder}[1]{\gdef\grothFolder{#1}}
\newcommand{\batch}[1]{\gdef\grothBatch{#1}}
\newcommand{\leaves}[2]{\gdef\grothFirst{#1}\gdef\grothLast{#2}}
\newcommand{\foldertitle}[1]{\gdef\grothFoldertitle{#1}}
\gdef\grothFolder{?}\gdef\grothBatch{?}\gdef\grothFirst{?}\gdef\grothLast{?}\gdef\grothFoldertitle{}

% --- Le résumé --------------------------------------------------------------
%
% Il ouvre la lecture modernisée, sous le titre « Résumé », et il est composé
% en petit. Ce n'est pas un ornement : c'est le seul passage du document qui
% ne soit pas des mathématiques, et un lecteur doit pouvoir le reconnaître
% avant de l'avoir lu — puis le sauter s'il connaît déjà le sujet, ou s'y
% arrêter s'il ne le connaît pas. Le filet qui le suit marque la reprise du
% corps.

\newenvironment{resume}
  {\small\setlength{\parskip}{0.45em}}
  {\par\medskip\hrule\medskip}

% --- Mots-clés ---------------------------------------------------------------
%
% En anglais, en clôture du résumé de la lecture modernisée : le vocabulaire
% moderne sous lequel chercher ce que les feuillets construisent. Le manifeste
% du site (`npm run manifest`) les extrait pour étiqueter la cote entière —
% cette ligne est la source unique des tags, il n'y a pas de fichier à côté.
\newcommand{\keywords}[1]{\par\smallskip\noindent
  {\footnotesize\textsc{Keywords} --- \textit{#1}}\par}

% --- L'appareil critique ----------------------------------------------------

% Le numéro de feuillet, en marge. C'est l'ancre qui permet de régler un
% désaccord sur une formule en regardant *un* feuillet plutôt qu'une chemise
% de six cents. Le site s'en sert aussi pour tourner les pages du fac-similé
% quand on fait défiler la transcription.
\newcommand{\leaf}[1]{%
  \marginnote{\footnotesize\color{gray!70}\textsf{#1}}%
  \label{leaf:#1}\ignorespaces}

% Illisible : on ne devine pas. Un mot inventé qui se lit comme les autres est
% le pire résultat possible.
\newcommand{\ill}{\textcolor{red!60!black}{[\,\ldots\,]}}

% Lecture douteuse : le mot est donné, et signalé comme incertain.
\newcommand{\uncertain}[1]{\uline{#1}}

% Ajout de l'éditeur — une lettre manquante, un mot sous-entendu.
\newcommand{\add}[1]{\textcolor{blue!50!black}{[#1]}}

% Biffé par Grothendieck lui-même. Ce qu'il a rayé fait partie du document :
% une rature dit souvent où la pensée a changé de direction.
\newcommand{\struck}[1]{\sout{#1}}

% Note du transcripteur, sur l'état matériel du feuillet.
\newcommand{\note}[1]{\par\smallskip{\small\color{gray!60!black}\textit{[#1]}}\par\smallskip}

% Note marginale de Grothendieck, distincte du corps du texte.
\newcommand{\marginal}[1]{\par\smallskip\noindent\hspace{1em}%
  {\small\color{gray!60!black}#1}\par\smallskip}

% --- Titre ------------------------------------------------------------------

\AtBeginDocument{%
  \begin{center}
    {\large\bfseries Fonds Alexandre Grothendieck --- cote n\textsuperscript{o}~\grothFolder}\\[2pt]
    {\itshape\grothFoldertitle}\\[4pt]
    {\small feuillets \grothFirst--\grothLast\ (lot \grothBatch)}\\[2pt]
    {\footnotesize\color{gray!60!black}Transcription établie d'après le fac-similé numérisé
     par l'Université de Montpellier.\\
     Les lectures incertaines sont soulignées, les illisibles notées [\,\ldots\,],
     les ajouts entre crochets.}
  \end{center}
  \vspace{1em}}

\endinput


\folder{81}
\pages{1}{54}
\dating{[à partir de 1981]}
\watermark{Édition de démonstration\\interprétation personnelle de l'œuvre}
\foldertitle{Immersions du disque et de la circonférence : notes manuscrites (s.d.) — lecture modernisée du dossier entier}

\begin{document}

\section*{Résumé}

\begin{resume}
Ces feuillets découpent. Un disque de papier qu'on entaille le long de
cordes, une sphère sur laquelle on a tracé un réseau de lignes, des points
rangés en cercle qu'on répartit en paquets : chaque fois la question est de
savoir comment les morceaux s'agencent, et ce qui, dans cet agencement, ne
dépend que de la combinatoire. Le titre que l'inventaire donne au dossier,
« Immersions du disque et de la circonférence », ne se lit sur aucun
feuillet et ne décrit aucun d'eux ; trois chemises de sa main portent
d'autres titres, qu'on suit ici.

Le dossier s'ouvre sur une remarque de quelques lignes, qui porte sur des
objets aujourd'hui bien connus sous un autre nom. Tracez sur une sphère un
graphe dont les arêtes ne se croisent pas : c'est une \emph{carte}. On peut
alors faire de la sphère une surface de Riemann et construire une fraction
rationnelle qui envoie les sommets sur $0$, les milieux d'arêtes --- et
rien d'autre --- sur $1$, les faces sur l'infini, et n'a pas d'autre valeur
critique que ces trois-là ; une fois trois points marqués, elle est unique.
Réciproquement, une fraction qui n'a pas d'autre valeur critique, et qui en
chaque point envoyé sur $1$ se comporte comme $z \mapsto z^2$ en $0$,
redessine une carte. Grothendieck affirme que les coefficients de ces
fractions sont des nombres algébriques, et donc que le groupe de Galois des
nombres algébriques agit sur les cartes, objets pourtant purement
topologiques. La page s'arrête au milieu d'une phrase. Ces cartes sont les
objets que l'\emph{Esquisse d'un programme} (1984) appelle \emph{dessins
d'enfants}, et ces fractions ce qu'on appelle des fonctions de Belyi ; les
feuillets ne citent ni l'une ni l'autre.

La partie la plus longue pose une question plus élémentaire. On coupe un
espace d'un seul tenant --- pensez à un disque --- par des coupures qui ne se
rencontrent pas, chacune le séparant en deux. Combien de morceaux, et
comment se touchent-ils ? Avec $n$ coupures on obtient $n+1$ morceaux, et si
l'on dessine un point par morceau et un trait par coupure, joignant les deux
morceaux qu'elle sépare, on obtient un \emph{arbre} : un réseau sans boucle,
comme le plan d'une maison où chaque porte est le seul passage d'une aile à
l'autre. Grothendieck écrit cette théorie trois fois. La rédaction qu'il
intitule « Version I » travaille avec les coupures elles-mêmes. Les deux
autres --- quelques pages sans titre, et la « Version II » --- remplacent
chaque coupure par ce qu'elle produit, deux « hémisphères » fermés
qui recouvrent l'espace et se rencontrent le long d'un « équateur », et
appellent \emph{non parallèles} deux tels découpages dont on peut choisir un
hémisphère de chacun sans point commun --- deux cordes d'un disque qui ne se
croisent pas. La Version II code tout cela par un ensemble fini
ordonné muni d'une symétrie (retourner une coupure), et en cherche les
axiomes. La proposition sur laquelle elle s'achève est fausse telle
qu'énoncée : trois cordes entourant un triangle la mettent en défaut, et la
démonstration s'interrompt.

La troisième partie pose la même question sur les seuls points d'un cercle.
On range $N$ points sur un cercle et on les répartit en paquets de sorte que
deux paquets ne « se croisent » jamais, c'est-à-dire que les polygones qu'ils
dessinent soient disjoints --- comme des convives autour d'une table ronde qui
se serreraient la main par groupes sans qu'aucun bras n'en croise un autre.
Ce sont les \emph{partitions non croisées}, qu'on étudie aujourd'hui sous ce
nom. Grothendieck montre qu'il y a toujours au moins deux paquets faits de
points consécutifs, sauf s'il n'y a qu'un paquet, et commence à les retirer
un à un, comme on effeuille un arbre par ses feuilles. Il décrit aussi ces
partitions par des familles d'arcs emboîtés, et remarque, sur trois cordes
parallèles, que la famille ne se retrouve pas à partir de la partition. Une
dernière page revient aux cartes : recoller deux cartes le long d'une arête,
et compter ce que deviennent les faces et le genre.

Ces pages montrent une manière de faire. Les définitions y portent leurs
corrections : « si elles sont distinctes » est ajouté en interligne, et la
démonstration en a besoin ; « strict.\ » se place devant « non
parallèles » dans la marge d'une page, et la notion n'est définie que dix
pages plus loin. Ce qui n'est pas encore compris est étiqueté comme tel ---
« Structure à étudier », « Intuition », « il faut vérifier que c'est
transitif » --- plutôt que glissé sous une formule. Et plutôt qu'une
rédaction raccommodée, on trouve des rédactions numérotées, « Version I »,
« Version II », dont chacune part des définitions.

Les noms sous lesquels chercher la suite : dessins d'enfants et théorème de
Belyi ; arbre dual d'un système de coupures, scissions compatibles et
théorème de Buneman, « pocsets » et arbres de Dunwoody ; partitions non
croisées de Kreweras, familles laminaires, laminations du disque.

\keywords{dessin d'enfant, Belyi map, Riemann existence theorem, Galois action on dessins, chord diagram, dual tree, nested cuts, compatible splits, Buneman's splits-equivalence theorem, pocset, Dunwoody's tree construction, non-crossing partition, laminar family, convex position, gluing of maps, genus of a map}
\end{resume}

\pagerange{1}{54}
\section*{Le fil du dossier, et les conventions}

\subsection*{Les stations}

Le dossier n'est pas daté de sa main ; l'inventaire le date « à partir de
1981 » et le range dans le groupe « Géométrie et topologie combinatoire »,
avec les dossiers 69 à 89\footnote{Le titre de l'inventaire, « Immersions du
disque et de la circonférence », ne figure sur aucun des feuillets
transcrits, et aucun ne traite d'immersion. Les trois chemises (pages 1, 13
et 40) portent leurs propres titres, qui sont repris ici comme titres de
sections. On ne sait pas d'où vient celui de l'inventaire. Si c'était le
dessein de Grothendieck, on remarquera --- c'est une conjecture de notre part,
que rien sur les pages ne soutient --- que la question de savoir quelles
courbes fermées immergées dans le plan bordent un disque immergé a été
résolue en 1967 par S.~Blank en découpant le disque le long de cordes
(exposé de V.~Poénaru au Séminaire Bourbaki, 1967-1968) ; les disques coupés
de cordes des pages 9 à 11 et les découpages des pages 13 à 39 seraient
alors des outils pour cette question.}.

\begin{enumerate}
  \item \emph{Pages 1 à 5}, « Cartes sphériques en général » : les cartes
  sur la sphère munies d'un repère, demi-arêtes libres admises,
  correspondent aux fractions rationnelles normalisées en $0$, $1$,
  $\infty$, non ramifiées hors de ces trois valeurs et ramifiées au plus à
  l'ordre $2$ au-dessus de $1$ ; les cartes sans demi-arête libre, à celles
  dont la ramification au-dessus de $1$ est partout exactement $2$ (page 2) ; leurs
  coefficients sont algébriques et $\mathrm{Aut}(\mathbb{C})$ agit à travers
  $\mathrm{Gal}(\overline{\mathbb{Q}}/\mathbb{Q})$ (page 4). Les pages 3 et 5
  sont des brouillons barrés.
  \item \emph{Pages 9 à 11}, des dessins sans texte : disques coupés de
  cordes, chacun avec un petit graphe ; la page 10 porte en outre quelques
  formules sur les ensembles de parties.
  \item \emph{Pages 13 à 24}, « Décompositions $n$-aires (et découpages)
  d'espaces topologiques », \emph{Version I} : un espace connexe coupé par
  des fermés disjoints $E_i$, chacun de complémentaire à deux composantes ;
  deux coupures font cinq pièces, $n$ coupures font $n+1$ morceaux, chacun
  s'écrit de façon unique comme le complémentaire d'hémisphères fermés
  disjoints, et morceaux et coupures forment un arbre.
  \item \emph{Pages 25 à 29}, un autre départ : les « dispositions
  binaires » d'un ensemble, la non-parallélité, puis la construction
  inverse de la précédente, qui d'un arbre tire une famille de
  décompositions deux à deux non parallèles.
  \item \emph{Pages 31 à 39}, « Décompositions $n$-aires d'un espace
  topologique », \emph{Version II}, paragraphes 1.1 à 1.3 et formules (1) à
  (27) : décompositions binaires fermées, paires non parallèles et le
  diagramme (14), paires strictement non parallèles, l'ensemble ordonné $C$
  à involution d'une famille, le « quartette ordonné » et une proposition
  dont la démonstration s'interrompt.
  \item \emph{Pages 40 à 52}, « Relations d'équivalence non parallèles » :
  polygones combinatoires et leurs segments, parties non parallèles,
  partitions non parallèles --- les partitions non croisées ---, familles
  emboîtées de segments et leurs cœurs, « anglets », et les mêmes notions
  pour un « quasi-segment ».
  \item \emph{Pages 53 et 54}, « Attachement de deux arêtes » : recoller deux
  cartes le long d'une arête, et le bilan des faces et du genre.
\end{enumerate}

\subsection*{Trois rédactions, gardées distinctes}

La théorie des découpages est écrite trois fois, et les trois ne disent pas
la même chose. La \emph{Version I} (pages 14 à 24) part des coupures $E_i$ et
des \emph{composantes ouvertes} $\Sigma_i^{\pm}$ de $X \smallsetminus E_i$ ;
elle ferme ces composantes en $\widehat{\Sigma}_i^{\pm} = \Sigma_i^{\pm} \cup
E_i$. La \emph{disposition binaire} des pages 25 à 29 est une paire de
parties d'un ensemble dont aucune ne contient l'autre. La \emph{Version II}
(pages 31 à 39) part de deux \emph{hémisphères fermés} $\Sigma^{\varepsilon}$
de réunion $X$, admet que l'un contienne l'autre, et n'appelle « propre » la
décomposition que si ce n'est pas le cas. Les lettres ne veulent donc pas
dire la même chose d'une rédaction à l'autre :
\[
\widehat{\Sigma} \ (\text{Version I}) = \Sigma \ (\text{Version II}), \qquad
\Sigma \ (\text{Version I}) = \Sigma^{0} \ (\text{Version II}).
\]
On garde, dans chaque section, les lettres de la rédaction qu'elle lit.

\subsection*{Conventions}

\begin{itemize}
  \item \emph{Non parallèle.} C'est le mot de Grothendieck, pour les
  décompositions (pages 25 et 32) comme pour les parties d'un polygone (pages
  43 et 51), et on le garde. Il veut dire « qui ne se croisent pas » : deux
  cordes disjointes d'un disque, que l'usage dirait parallèles, définissent
  des décompositions \emph{non parallèles} au sens de ces pages. Le nom
  actuel est \emph{compatibles} (pour des scissions) ou \emph{non croisées}
  (pour des parties d'un cycle).
  \item Un \emph{morceau} est une classe de la partition $\Pi_I^0$ du
  complémentaire des coupures (la page dit « classe »). Le graphe des
  morceaux de la page 23, que la page note $(S, I, R)$, est noté ici
  $\Gamma_I$.
  \item Un \emph{polygone} est un ensemble fini $S$ de cardinal $N \geqslant
  3$ muni d'une « bi-circulation » $\{u, u^{-1}\}$, où $u$ est une permutation
  circulaire de $S$ ; un \emph{segment} est une partie formée de sommets
  consécutifs, non vide et distincte de $S$ ; $\partial T$ est l'ensemble de
  ses extrémités. La relation d'équivalence des pages 47 à 50 est notée $R$,
  ses classes $S_i$ ; l'ensemble résiduel de la page 46, que la page note
  aussi $R$, est noté $Z$.
\end{itemize}

Restent des collisions de lettres, qu'on ne corrige pas parce que les
sections ne se mêlent pas : $S$ est l'ensemble des morceaux (page 23), celui
des sommets d'un arbre (page 27) et le polygone (page 41) ; $I$ indexe les
coupures (pages 14 à 24) et désigne le quasi-segment (page 51) ; $A$ indexe
une famille (page 37), est l'ensemble des arêtes d'un arbre (page 27) et une
partie du polygone (pages 42 et suivantes).

\subsection*{Ce que seule une lecture d'ensemble peut dire}

Quatre liens, qu'aucun feuillet n'écrit, et qui sont de nous. La Version I
va des coupures à l'arbre (pages 16 à 24) ; les pages 27 à 29 vont de l'arbre
aux décompositions : ce sont les deux sens d'une même correspondance, que la
Version II cherche à axiomatiser. Le contre-exemple à la proposition de la
page 39 est la configuration de trois cordes autour d'un triangle, dont le
graphe dessiné page 11 est un tripode. Les trois cordes parallèles de la page
47, qui montrent qu'une famille de segments ne se retrouve pas à partir de
sa partition, rappellent la rangée de quatre cordes parallèles marquée
« ? » page 9. Et la bi-circulation de la page 41 est exactement la troisième
description d'un polygone combinatoire du dossier 89 (pages 19 et 20) : une
paire de permutations circulaires inverses l'une de l'autre.

\subsection*{Ce que le dossier annonce sans l'établir}

La fin de la page 4 (« le diviseur de $f$ est »). La question, s'il y en a
une, des dessins des pages 9 à 11. L'égalité annoncée page 16, le lemme de la
page 22 et le lemme de l'arbre de la page 23, dont la démonstration commence
page 24 et s'arrête ; on les démontre ici. La « description complète » de la
page 28 et le cas des ensembles vides, « à examiner de plus près » page 29.
La transitivité demandée en marge de la page 38 ; la réciproque qu'annonce
l'« Inversement » de la même page. La proposition de la page 39, fausse telle
qu'énoncée. Le « plus petit » couple de la page 43. L'« Intuition » de la
page 44. La reconstruction d'une famille de segments à partir de sa partition
(page 47), que le NB de la même page réfute. La tour de la page 50. Le NB et
l'étude de $\Sigma_{\max}$ de la page 52. La dérivation des bilans de la page
53.

\pagerange{1}{5}
\section*{I. Cartes sphériques en général (pages 1 à 5)}

\pagerange{2}{2}
\subsection*{Le dictionnaire (page 2)}

Une \emph{carte sphérique orientée} est un graphe fini tracé sans croisement
sur une sphère orientée, dont le complémentaire est une réunion de disques
ouverts, les \emph{faces} ; on la considère à isotopie près. Une carte
\emph{à demi-arêtes libres} admet en outre des \emph{demi-arêtes libres},
arcs attachés à un sommet par un bout et libres de l'autre, son
complémentaire étant toujours une réunion de disques ouverts ; les cartes
sont celles qui n'en ont pas. Un \emph{repère}
y fixe trois points\footnote{La page n'explique pas le mot « repère ». La
normalisation $f(0)=0$, $f(1)=1$, $f(\infty)=\infty$ qui suit revient à
marquer un sommet (au-dessus de $0$), un milieu d'arête ou l'extrémité libre
d'une demi-arête (au-dessus de $1$) et une face (au-dessus de $\infty$) :
c'est notre lecture. La page ne dit pas s'ils doivent être incidents ; la
normalisation ne le demande pas.}. La page énonce trois descriptions
équivalentes de ces objets.

\textbf{Premier dictionnaire.} Il appelle \emph{spéciales} les fractions
rationnelles $f \colon \mathbb{P}^1_{\mathbb{C}} \to \mathbb{P}^1_{\mathbb{C}}$
telles que
\[
f(0) = 0, \qquad f(1) = 1, \qquad f(\infty) = \infty,
\]
$f$ est étale au-dessus de $\mathbb{P}^1 \smallsetminus \{0, 1, \infty\}$, et
tout point critique $z \in \mathbb{C}$ de $f$ avec $f(z) \notin \{0,\infty\}$
vérifie $f(z) = 1$ et $f''(z) \neq 0$ : au-dessus de $1$, l'indice de
ramification est au plus $2$\footnote{La page définit d'abord « spécial »
par les trois valeurs et l'étalité hors de $\{0,1,\infty\}$, et n'ajoute
qu'à la ligne suivante, « et de plus », la condition $f''(z) \neq 0$. Sans
elle, la première équivalence est fausse : les fonctions étales hors de
trois points correspondent aux \emph{hypercartes} (cartes biparties), et
c'est la ramification au plus $2$ au-dessus de $1$ qui fait de chaque point
au-dessus de $1$ un milieu d'arête ou l'extrémité libre d'une demi-arête.
On a intégré la condition à la définition.}. Alors :
\begin{itemize}
  \item les fonctions spéciales correspondent aux cartes sphériques
  orientées repérées \emph{à demi-arêtes libres} ;
  \item parmi elles, celles dont \emph{tout} point au-dessus de $1$ est
  d'indice exactement $2$ correspondent aux cartes sphériques orientées
  repérées, sans demi-arête libre.
\end{itemize}
Les sommets sont les points au-dessus de $0$, avec pour degré l'indice de
ramification ; les faces, les points au-dessus de $\infty$ ; les arêtes, les
points au-dessus de $1$ d'indice $2$, et les demi-arêtes libres, ceux
d'indice $1$ --- l'arête ou la demi-arête étant l'adhérence de la composante
de $f^{-1}(]0,1])$ qui passe par ce point\footnote{La page n'énonce
que l'équivalence, entre « carte sphérique » et fonction spéciale, et ne
parle ni d'arêtes ni de demi-arêtes. Sa condition est « au plus $2$ » : elle
ne demande $f''(z) \neq 0$ qu'aux points critiques, et admet donc au-dessus
de $1$ des points d'indice $1$, extrémités libres de demi-arêtes. Prise au
sens usuel, sans demi-arête libre, une carte ne correspond qu'aux fonctions
d'indice exactement $2$ en tout point au-dessus de $1$ (fonctions de Belyi
dites « propres », en anglais \emph{clean}) ; oublier les demi-arêtes ne
donne pas une bijection, puisque $T$ et $T^2$ ont alors le même graphe, un
sommet sans arête. Le dictionnaire est de nous.}.

C'est une forme du \emph{théorème d'existence de Riemann}. Une carte fait de
la sphère privée des sommets, des milieux d'arêtes, des extrémités libres
des demi-arêtes et des centres de faces un revêtement fini de la sphère privée de trois points ; ce revêtement porte une
unique structure complexe qui rende la projection holomorphe, elle se
prolonge aux points ôtés, et la surface obtenue, de genre $0$, est isomorphe à
$\mathbb{P}^1$, l'isomorphisme étant fixé par les trois points du
repère\footnote{Ces trois étapes (existence de Riemann, prolongement, $\mathbb{P}^1$
comme seule surface de Riemann compacte de genre $0$, rigidité d'un
automorphisme qui fixe trois points) sont sous-entendues par la page. Le
repère est ce qui rend les objets rigides : une carte repérée n'a pas
d'automorphisme non trivial.}. On reconnaît ce qu'on appelle aujourd'hui une
\emph{fonction de Belyi} sur $\mathbb{P}^1$, et la carte est son \emph{dessin
d'enfant}\footnote{Aucun des deux noms n'est sur la page : « dessin
d'enfant » est le terme de l'\emph{Esquisse d'un programme} (1984), et
« fonction de Belyi » vient du théorème de G.~V.~Belyi (1979), selon lequel
une courbe sur $\mathbb{C}$ est définie sur $\overline{\mathbb{Q}}$ si et
seulement si elle admet une telle fonction. Ici la courbe est
$\mathbb{P}^1$, et le théorème de Belyi n'intervient pas. La date de
l'inventaire rend possible que Grothendieck ait connu le travail de Belyi ;
rien sur la page ne le cite. Voir aussi la fin du dossier 88 (page 15), qui
s'arrête au bord de la même observation pour ses cartes de type $(p,q)$.}.

\textbf{Deuxième dictionnaire.} Écrivons $f = P/Q$ avec $P, Q \in
\mathbb{C}[z]$ premiers entre eux. Les conditions aux trois points
s'écrivent
\[
P(0) = 0, \qquad P(1) = Q(1), \qquad \deg Q < \deg P,
\]
d'où $Q(0) \neq 0$ et $P(1) = Q(1) \neq 0$ par primalité ; on normalise par
$P(1) = Q(1) = 1$. En un point $z$ où $P(z) Q(z) \neq 0$, on a
$f'/f = P'/P - Q'/Q$, et la condition sur les points critiques devient :
\[
\frac{P'(z)}{P(z)} = \frac{Q'(z)}{Q(z)} \ \Longrightarrow\ P(z) = Q(z)
\ \text{ et } \ P''(z) \neq Q''(z).
\]
La seconde conclusion vient de ce qu'en un tel point
$f''(z) = \bigl(P''(z) - Q''(z)\bigr)/Q(z)$, puisque $P(z) = Q(z)$ et
$P'(z) = Q'(z)$\footnote{La page écrit $\frac{P''(z)}{P(z)} -
\frac{Q''(z)}{Q(z)} \neq 0$ « i.e. » $P''(z) \neq Q''(z)$, ce qui est
juste une fois $P(z) = Q(z)$ acquis. Le calcul de $(P/Q)''$ qui le justifie
est au verso, page 5.}.

\pagerange{2}{2}
\subsection*{Une question, et sa réponse (page 2)}

La page note l'identité
\[
Q P' - P Q' = (Q - P)\,P' + P\,(P' - Q').
\]
En un point critique $z$ de $f$ (où $QP' - PQ' = 0$) avec $P(z) \neq 0$, elle
montre que $P(z) = Q(z)$ entraîne $P'(z) = Q'(z)$, et que réciproquement
$P'(z) = Q'(z)$ entraîne $P(z) = Q(z)$ \emph{pourvu que} $P'(z) \neq 0$. D'où
la question de la page : peut-on avoir $P'(z) = Q'(z) = 0$ sans $P(z) =
Q(z)$ ?

La réponse est oui, et l'exemple que la page commence la
donne\footnote{L'exemple est au bas de la page, déchirée, et se poursuit dans
la marge gauche, lue en partie seulement ; la conclusion qui suit est de
nous.}. Prenons $Q = 1$ : les conditions se réduisent à $P(0) = 0$, $P(1) =
1$, et la question devient de savoir si $P'(z) = 0$ entraîne $P(z) \in \{0,
1\}$. Pour $P(T) = aT^2 + bT$ avec $a + b = 1$ et $a \neq 0$, le seul point
critique est $z = -b/2a$, et
\[
P\Bigl(-\frac{b}{2a}\Bigr) = -\frac{b^2}{4a}.
\]
Cette valeur vaut $0$ si et seulement si $b = 0$, et $1$ si et seulement si
$b^2 = -4a = 4b - 4$, c'est-à-dire $b = 2$. Dans cette famille, les seules
fonctions spéciales sont donc $T^2$ --- un sommet de degré $2$ et deux
demi-arêtes libres --- et $2T - T^2 = 1 - (1-T)^2$ --- deux sommets, en $0$ et
en $2$, joints par une arête dont le milieu est $1$ ; pour toute autre valeur
de $b$, $P'$ s'annule en un point où $P$ ne vaut ni $0$ ni $1$.

\pagerange{4}{4}
\subsection*{L'action de $\mathrm{Aut}(\mathbb{C})$ (page 4)}

Si $\sigma$ est un automorphisme du corps $\mathbb{C}$ et $f$ une fonction
spéciale, la fonction $f^{\sigma}$ obtenue en appliquant $\sigma$ aux
coefficients est encore spéciale : toutes les conditions sont algébriques.
Donc $\mathrm{Aut}(\mathbb{C})$ agit sur l'ensemble des fonctions spéciales,
et, dit la page, il agit à travers
$\mathrm{Gal}(\overline{\mathbb{Q}}/\mathbb{Q})$ parce que les coefficients
de $P$ et $Q$ sont algébriques ; il revient au même de dire que les zéros de
$P$ et de $Q$ le sont.

La page affirme l'algébricité sans la justifier\footnote{La page s'arrête sur
« i.e. le diviseur de $f$ est », au milieu du feuillet. L'argument qui suit
est de nous.}. Voici pourquoi elle est vraie. Pour un degré $d$ donné, il n'y
a qu'un nombre fini de fonctions spéciales : un revêtement de degré $d$ de
$\mathbb{P}^1 \smallsetminus \{0,1,\infty\}$ est donné par une action
transitive du groupe libre à deux générateurs sur un ensemble à $d$
éléments, il y en a un nombre fini à isomorphisme près, et chacun n'admet
qu'un nombre fini de repères. Normalisés par $P(1) = Q(1) = 1$, les
coefficients sont déterminés par $f$ ; chacun n'a donc qu'un nombre fini de
conjugués sous $\mathrm{Aut}(\mathbb{C})$, et un nombre complexe qui n'a
qu'un nombre fini de conjugués est algébrique, puisque $\mathrm{Aut}(\mathbb{C})$
agit transitivement sur les nombres transcendants. L'action se factorise
alors par la restriction $\mathrm{Aut}(\mathbb{C}) \to
\mathrm{Gal}(\overline{\mathbb{Q}}/\mathbb{Q})$.

On a ainsi un groupe de Galois qui agit sur des objets purement
combinatoires --- des graphes tracés sur une sphère. C'est l'observation que
l'\emph{Esquisse d'un programme} place au centre, pour les cartes de tout
genre, et dont elle tire, à l'aide du théorème de Belyi, que l'action sur
l'ensemble de tous les dessins est fidèle\footnote{La fidélité n'est pas sur
la page, qui ne considère que le genre $0$.}.

\pagerange{3}{5}
\subsection*{Les versos (pages 3 et 5)}

Les pages 3 et 5, barrées, sont des brouillons. La page 3 calcule dans une
algèbre engendrée par des éléments $e_i$, $f_i$ et $\eta$, munie d'une
dérivation $\delta$, et conclut, sous ses hypothèses, que $\delta \eta = \pm
e_{d-2}$ ; rien ne la rattache aux cartes\footnote{Transcription condensée,
formules seulement ; on ne cherche pas à reconstituer la structure étudiée.}.
La page 5 porte, sur une couche, la suite de ce calcul, et, tête-bêche, la
dérivée seconde
\[
\Bigl(\frac{P}{Q}\Bigr)'' = \frac{P''Q^2 - Q''PQ - 2P'QQ' + 2PQ'^2}{Q^3},
\]
qui sert à la page 2.

\pagerange{9}{11}
\section*{II. Disques coupés de cordes (pages 9 à 11)}

Les pages 9 et 11 sont des pages de dessins, sans texte. Chacune montre des
rangées de disques coupés par des cordes, droites ou courbes, parfois
concourantes, une région parfois hachurée ; sous chaque disque, un petit
graphe. Pour la plupart des dessins, les sommets du graphe semblent être les
cordes et ses arêtes leurs croisements : un point pour une corde, deux points
pour deux cordes disjointes, un segment pour deux cordes qui se croisent, un
triangle pour trois cordes deux à deux sécantes\footnote{Cette lecture est
celle du transcripteur (« apparemment ») ; un dessin au moins, trois cordes
concourantes au-dessus d'un segment, ne s'y conforme pas, et on ne force
pas l'interprétation. Le graphe des croisements d'une famille de cordes
s'appelle aujourd'hui un \emph{graphe de cercle} ; rien sur la page ne dit
que c'est ce qu'il étudie.}. Plusieurs dessins sont marqués « ok », avec
des variantes introduites par « ou » ; sous une rangée de disques à quatre
cordes parallèles, un « ? ». La page 11 ajoute, parmi les graphes, un
tripode en forme de Y, un carré, un triangle surmonté d'un segment.

Ces dessins sont des \emph{diagrammes de cordes}, et ils reviennent plus
loin sous deux formes : les partitions en classes de deux éléments de la
page 44, qui sont des familles de cordes sans croisement, et les trois
cordes parallèles de la page 47. Les configurations à cordes sécantes que
dessinent les pages 9 et 11 sont exactement celles que la condition « non
parallèle » des pages 40 à 52 exclut ; c'est un rapprochement de notre
part.

\textbf{La page 10.} Au milieu de ces dessins, quelques formules sur les
ensembles de parties $\mathbb{P}(X)$ et les applications :
\[
\mathbb{P}(X \times Z) \simeq \mathrm{Hom}(Z \times X, \mathcal{E}) =
\mathrm{Hom}\bigl(X, \mathrm{Hom}(Z, \mathcal{E})\bigr), \qquad
\mathrm{Hom}(X, Y) \subset \mathbb{P}(X \times Y).
\]
Si l'on lit $\mathcal{E}$, que la page ne définit pas, comme un ensemble à
deux éléments --- ou, dans un topos, le classifiant des sous-objets ---, la
première ligne est la correspondance entre parties et fonctions
caractéristiques, suivie de la curryfication, et vaut sans le point
d'interrogation que la page met sur « $\simeq$ » ; la seconde identifie une
application à son graphe\footnote{Une autre ligne porte
$\mathrm{Hom}(\mathcal{E}, \mathrm{Hom}(X,Z)) =
\mathbb{P}(\mathrm{Hom}(X,Z))$, dont la première lettre est douteuse. Telle
qu'écrite, la variance est fausse : $\mathbb{P}(Y) = \mathrm{Hom}(Y,
\mathcal{E})$ est contravariant en $Y$, et il faudrait
$\mathrm{Hom}(\mathrm{Hom}(X,Z), \mathcal{E})$. La lecture de $\mathcal{E}$
comme classifiant est la nôtre.}. Le croquis voisin, $H \subset P \times X
\times Y$ au-dessus de $P \times X$, puis de $P$, évoque la relation
universelle sur un objet de parties, par laquelle on construit les
exponentielles d'un topos à partir de ses objets de parties ; la page n'en
dit pas plus. Rien ne rattache cette page au reste du dossier.

\pagerange{13}{24}
\section*{III. Découpages d'un espace par des coupures : la Version I (pages 13 à 24)}

\pagerange{13}{14}
\subsection*{Les hypothèses (pages 13 et 14)}

Soient $X$ un espace topologique \emph{connexe} et $(E_i)_{i \in I}$ une
famille \emph{finie} de parties fermées de $X$ telles que :
\begin{enumerate}
  \item[(a)] pour tout $i$, $X \smallsetminus E_i$ a exactement deux
  composantes connexes, notées $\Sigma_i^{+}$ et $\Sigma_i^{-}$ ;
  \item[(b)] pour $i \neq j$, $E_j$ est contenu dans l'une des deux
  composantes de $X \smallsetminus E_i$.
\end{enumerate}
La condition (b) entraîne que les $E_i$ sont deux à deux disjoints ; elle
est automatique si les $E_i$ sont connexes et disjoints. On pose
$\widehat{\Sigma}_i^{\pm} = \Sigma_i^{\pm} \cup E_i$\footnote{La connexité de
$X$ et la finitude de $I$ ne sont pas dans l'énoncé. La première apparaît en
marge (« ce qui est le cas si $X$ est connexe ») et sert à chaque pas ; la
seconde est demandée par la récurrence sur $\operatorname{card} I$ et par le
fait que les morceaux soient ouverts. Les deux sont supposées dans toute la
section.}.

Trois conséquences servent partout. Les $\Sigma_i^{\pm}$ sont \emph{ouverts}
dans $X$ : chacun est fermé dans $X \smallsetminus E_i$ comme composante, donc
ouvert comme complémentaire de l'autre, et $X \smallsetminus E_i$ est ouvert
dans $X$\footnote{La marge dit « ouvertes », lecture incertaine. C'est vrai
sans hypothèse de connexité locale, parce qu'il n'y a que deux
composantes.}. On a $\overline{\Sigma_i^{\pm}} \subset \Sigma_i^{\pm} \cup
E_i$, puisque l'autre composante est un ouvert disjoint ; en particulier
$\widehat{\Sigma}_i^{\pm}$ est fermé. Enfin, $X$ étant connexe,
$\overline{\Sigma_i^{\pm}} \cap E_i \neq \emptyset$ : sinon $\Sigma_i^{\pm}$
serait ouvert et fermé, non vide et distinct de $X$. En particulier $E_i \neq
\emptyset$.

L'exemple à garder en tête est celui d'un disque fermé coupé par des cordes
disjointes. Pour des cordes rectilignes, (a) est élémentaire ; pour des arcs
simples joignant deux points du bord, c'est une forme du théorème de Jordan,
et il faut celui de Schoenflies pour savoir que les morceaux sont encore des
disques\footnote{La page ne mentionne ni l'un ni l'autre ; c'est ce qu'il
faut pour que l'exemple du disque, qu'elle dessine page 14, satisfasse ses
hypothèses.}.

\pagerange{14}{15}
\subsection*{Deux coupures (pages 14 et 15)}

Soit $I = \{1, 2\}$, et notons $\Sigma_i^{+}$ la composante de $X
\smallsetminus E_i$ qui \emph{ne contient pas} l'autre coupure.

\textbf{Les deux calottes sont disjointes.} On a
\[
\Sigma_1^{+} \subset \Sigma_2^{-}, \qquad \Sigma_2^{+} \subset \Sigma_1^{-},
\qquad \Sigma_1^{+} \cap \Sigma_2^{+} = \emptyset .
\]
En effet $\Sigma_2^{-}$ contient $E_1$, donc en est un voisinage ouvert, et
$\overline{\Sigma_1^{+}}$ rencontre $E_1$ ; donc $\Sigma_2^{-}$ rencontre
$\Sigma_1^{+}$. La réunion $\Sigma_2^{-} \cup \Sigma_1^{+}$ est alors un
connexe contenu dans $X \smallsetminus E_2$ (car $\Sigma_1^{+} \cap E_2 =
\emptyset$) et contenant la composante $\Sigma_2^{-}$ : elle lui est égale,
et $\Sigma_1^{+} \subset \Sigma_2^{-}$. Symétriquement pour l'autre
inclusion, et la disjonction en résulte.

\textbf{Les cinq pièces.} La partition commune des deux partitions $X =
\Sigma_i^{+} \sqcup E_i \sqcup \Sigma_i^{-}$ est
\[
X = \Sigma_1^{+} \sqcup E_1 \sqcup \Sigma_{12} \sqcup E_2 \sqcup
\Sigma_2^{+}, \qquad \Sigma_{12} = \Sigma_1^{-} \cap \Sigma_2^{-},
\]
avec $\Sigma_2^{-} = \Sigma_1^{+} \sqcup E_1 \sqcup \Sigma_{12}$ et
$\Sigma_1^{-} = \Sigma_{12} \sqcup E_2 \sqcup \Sigma_2^{+}$. Les trois pièces
$\Sigma_1^{+}$, $\Sigma_{12}$, $\Sigma_2^{+}$ sont ouvertes, et
\[
\Sigma_1^{+} \subsetneq \overline{\Sigma_1^{+}} \subset \Sigma_1^{+} \cup
E_1, \qquad \overline{\Sigma_{12}} \subset \Sigma_{12} \cup E_1 \cup E_2,
\qquad \Sigma_2^{+} \subsetneq \overline{\Sigma_2^{+}} \subset \Sigma_2^{+}
\cup E_2 .
\]
\textbf{La pièce du milieu touche les deux coupures} :
$\overline{\Sigma_{12}} \cap E_1 \neq \emptyset$ et $\overline{\Sigma_{12}}
\cap E_2 \neq \emptyset$ ; en particulier $\Sigma_{12} \neq \emptyset$. Si
l'on avait $\overline{\Sigma_{12}} \cap E_1 = \emptyset$, on aurait
$\overline{\Sigma_{12}} \subset \Sigma_{12} \cup E_2$, puis
\[
\overline{\Sigma_1^{-}} = \overline{\Sigma_2^{+}} \cup E_2 \cup
\overline{\Sigma_{12}} \subset \Sigma_1^{-},
\]
et $\Sigma_1^{-}$ serait fermé, ce que la connexité interdit.

\pagerange{16}{16}
\subsection*{Le cas des frontières communes (page 16)}

Supposons de plus que chaque coupure soit la frontière commune de ses deux
côtés :
\[
\overline{\Sigma_i^{+}} = \Sigma_i^{+} \cup E_i, \qquad
\overline{\Sigma_i^{-}} = \Sigma_i^{-} \cup E_i \qquad (i \in I),
\]
comme c'est le cas pour une corde d'un disque. Alors
\[
\overline{\Sigma_{12}} = \Sigma_{12} \cup E_1 \cup E_2 .
\]
La page l'annonce (« je dis que ») sans le démontrer\footnote{La page écrit
« les deux comp.\ conn.\ de $X \smallsetminus \Sigma_i$ » pour $X
\smallsetminus E_i$, ici comme page 17. La démonstration est de nous.}. En
effet $E_1 \subset \overline{\Sigma_1^{-}} = \overline{\Sigma_2^{+}} \cup E_2
\cup \overline{\Sigma_{12}}$, et $E_1$ ne rencontre ni $E_2$ ni
$\overline{\Sigma_2^{+}} \subset \Sigma_2^{+} \cup E_2$ ; donc $E_1 \subset
\overline{\Sigma_{12}}$, et de même pour $E_2$.

\pagerange{16}{19}
\subsection*{Compter les morceaux (pages 16 à 19)}

Pour $J \subset I$, notons $U_J = X \smallsetminus \bigcup_{j \in J} E_j$ et
$\Pi_J^0$ la partition de $U_J$ qui est la partition commune des partitions
$\{\Sigma_j^{+} \cap U_J, \Sigma_j^{-} \cap U_J\}$ ($j \in J$) : deux points
sont dans la même classe s'ils sont, pour chaque $j$, du même côté de $E_j$.
Ses classes non vides sont les \emph{morceaux}\footnote{Ce sont les
composantes connexes de $U_J$ dans l'exemple du disque, mais pas en général.
Sur un cercle, une paire de points $E = \{p, q\}$ est une coupure au sens de
(a) ; avec deux paires non entrelacées, dans l'ordre $p_1, q_1, p_2, q_2$, le
morceau du milieu $\Sigma_{12}$ est la réunion des deux arcs $]q_1, p_2[$ et
$]q_2, p_1[$. L'exemple est de nous ; le compte qui suit porte sur les
classes, et reste juste.}. Ils sont ouverts, comme intersections finies
d'ouverts.

\textbf{Le compte (pages 17 et 19).} \emph{Si $\operatorname{card} I = n$,
la partition $\Pi_I^0$ a exactement $n + 1$ morceaux.}

Par récurrence sur $n$, les cas $n = 0, 1$ étant clairs. Soit $I = J \sqcup
\{i\}$. Par (b), $E_i$ est, pour chaque $j \in J$, d'un seul côté de $E_j$ ;
il est donc contenu dans un unique morceau de $\Pi_J^0$, noté $U_{J,i}$, qui
en est un voisinage ouvert.

\textbf{Lemme (page 17).} \emph{Tout morceau de $\Pi_J^0$ distinct de
$U_{J,i}$ est contenu dans $\Sigma_i^{+}$ ou dans $\Sigma_i^{-}$.}

« C'est presque une tautologie. » Un morceau $U_\alpha \neq U_{J,i}$ est, pour
un certain $j \in J$, du côté de $E_j$ qui ne contient pas $E_i$ ; par
l'étude de deux coupures, ce côté est contenu dans le côté de $E_i$ qui
contient $E_j$, et $U_\alpha$ aussi.

Les morceaux de $\Pi_I^0$ sont donc ceux de $\Pi_J^0$ autres que $U_{J,i}$,
inchangés, et les deux ensembles $U_{J,i} \cap \Sigma_i^{+}$ et $U_{J,i} \cap
\Sigma_i^{-}$, qui sont non vides parce que $U_{J,i}$ est un voisinage de
$E_i$ et que les deux côtés de $E_i$ ont $E_i$ dans leur adhérence. Il y en
a $(n - 1) + 2 = n + 1$ (page 19).

\pagerange{18}{22}
\subsection*{La forme des morceaux (pages 18 à 22)}

\textbf{Lemme (page 19).} \emph{Tout morceau de $\Pi_I^0$ s'écrit}
\[
V = X \smallsetminus \bigcup_{\alpha \in K} \widehat{\Sigma}_\alpha^{+}
= \bigcap_{\alpha \in K} \Sigma_\alpha^{-},
\]
\emph{où $K \subset I$ est non vide si $I$ l'est, et où, pour chaque $\alpha
\in K$, $\Sigma_\alpha^{+}$ est l'une des deux composantes de $X
\smallsetminus E_\alpha$ (et $\Sigma_\alpha^{-}$ l'autre), les
$\widehat{\Sigma}_\alpha^{+}$ étant deux à deux disjoints.}

Pour $\alpha \neq \beta$, la disjonction de $\widehat{\Sigma}_\alpha^{+}$ et
$\widehat{\Sigma}_\beta^{+}$ équivaut à celle de $\Sigma_\alpha^{+}$ et
$\Sigma_\beta^{+}$ : par l'étude de deux coupures, une seule des quatre
paires de composantes est disjointe, celle des deux côtés « extérieurs », et
chaque coupure est alors du côté intérieur de l'autre.

La démonstration (pages 18 et 20) est une récurrence. Par hypothèse de
récurrence, $U_{J,i} = X \smallsetminus \bigcup_{\alpha \in K_0}
\widehat{\Sigma}_\alpha^{+}$, et
\[
U_{J,i} \cap \Sigma_i^{-} = X \smallsetminus \Bigl( \bigcup_{\alpha \in K_0}
\widehat{\Sigma}_\alpha^{+} \cup \widehat{\Sigma}_i^{+} \Bigr).
\]
Dans la réunion on peut omettre les $\alpha$ tels que $\Sigma_\alpha^{+}
\subset \Sigma_i^{+}$, puisqu'alors $\widehat{\Sigma}_\alpha^{+} \subset
\widehat{\Sigma}_i^{+}$ ; soit $K_1$ l'ensemble des autres. Reste à voir que
pour $\alpha_0 \in K_1$, $\Sigma_{\alpha_0}^{+} \cap \Sigma_i^{+} =
\emptyset$. Or $E_i \subset U_{J,i} \subset \Sigma_{\alpha_0}^{-}$, donc
$\Sigma_{\alpha_0}^{+}$ est le côté de $E_{\alpha_0}$ qui ne contient pas
$E_i$ ; il est contenu dans le côté de $E_i$ qui contient $E_{\alpha_0}$.
Comme $\Sigma_{\alpha_0}^{+} \not\subset \Sigma_i^{+}$, ce côté est
$\Sigma_i^{-}$, et $\Sigma_i^{+}$ est l'autre, disjoint de
$\Sigma_{\alpha_0}^{+}$ par l'étude de deux coupures\footnote{La page écrit
« $\Sigma_{\alpha_0} \cap \Sigma_i = \emptyset$ » sans les signes. Elle
annonce la conclusion (« ce qui achèvera la démonstration ») avant
l'argument, qui occupe le bas de la page 20 ; on a rétabli l'ordre.}. On
conclut avec $K = K_1 \cup \{i\}$ ; l'autre morceau, $U_{J,i} \cap
\Sigma_i^{+}$, se traite de même.

\textbf{Corollaire du lemme (page 20).} \emph{L'écriture de $V$ est unique :}
\[
K = \{ \alpha \in I \mid \overline{V} \cap E_\alpha \neq \emptyset \},
\]
\emph{et, pour $\alpha \in K$, $\Sigma_\alpha^{+}$ est la composante de $X
\smallsetminus E_\alpha$ qui ne rencontre pas $V$.} Comme $V$ est ouvert et
ne rencontre aucune coupure, la condition $\overline{V} \cap E_\alpha \neq
\emptyset$ équivaut à $\dot{V} \cap E_\alpha \neq \emptyset$, où $\dot{V}$
est la frontière de $V$.

Les pages 21 et 22 le démontrent en deux temps.

\emph{Si $\alpha \in K$, alors $\overline{V} \cap E_\alpha =
\overline{\Sigma_\alpha^{-}} \cap E_\alpha$, qui est non vide.} Posons $Y =
\widehat{\Sigma}_\alpha^{+}$ et $Z = \bigcup_{\beta \in K \smallsetminus
\{\alpha\}} \widehat{\Sigma}_\beta^{+}$, fermés disjoints. On a $X
\smallsetminus Y = (X \smallsetminus (Y \cup Z)) \cup Z$, donc, $Z$ étant
fermé, $\overline{X \smallsetminus Y} = \overline{X \smallsetminus (Y \cup
Z)} \cup Z$, et en coupant par $Y$, qui ne rencontre pas $Z$ :
\[
\overline{X \smallsetminus (Y \cup Z)} \cap Y = \overline{X \smallsetminus Y}
\cap Y ,
\]
ce qui est l'égalité cherchée (« immédiat », dit la page, après une première
tentative encadrée et barrée).

\emph{Si $\beta \notin K$, alors $\overline{V} \cap E_\beta = \emptyset$.} Le
morceau $V$ ne rencontre aucune coupure, donc $E_\beta \subset X
\smallsetminus V = \bigcup_{\alpha \in K} \widehat{\Sigma}_\alpha^{+}$, et,
les coupures étant disjointes, $E_\beta \subset \bigcup_{\alpha \in K}
\Sigma_\alpha^{+}$, ouvert disjoint de $V$.

\textbf{Lemme (page 22).} \emph{Mieux : il existe $\alpha \in K$ tel que
$E_\beta \subset \Sigma_\alpha^{+}$.} Par (b), $E_\beta$ est, pour chaque
$\alpha \in K$, contenu dans $\Sigma_\alpha^{+}$ ou dans $\Sigma_\alpha^{-}$ ;
s'il était contenu dans tous les $\Sigma_\alpha^{-}$, il le serait dans leur
intersection $V$, qu'il ne rencontre pas, et il est non vide\footnote{La page
énonce le lemme pour une famille quelconque de $\Sigma_\alpha^{+}$ deux à
deux disjoints et laisse la démonstration en suspens (« Cela provient
… ») ; la marge dit que c'est « trivial, du fait que les $E_i$ sont
disjoints ». La disjonction ne donne que l'inclusion dans la réunion ; pour
passer à un seul $\Sigma_\alpha^{+}$, il faut (b), ou la connexité de
$E_\beta$. La démonstration est de nous.}.

\textbf{Corollaire (page 22).} \emph{Les morceaux de la partition
$\Pi_K^0$, définie par les seules coupures $E_\alpha$ ($\alpha \in K$), sont
les $\Sigma_\alpha^{+}$ ($\alpha \in K$) et $V$.} Chaque $\Sigma_\alpha^{+}$
est en effet contenu dans tous les $\Sigma_\beta^{-}$ ($\beta \neq \alpha$),
et le compte, $\operatorname{card} K + 1$, est le bon (« immédiat par
récurrence »). L'unicité de l'écriture en résulte : $K$ est l'ensemble des
coupures qui touchent $V$, et $\Sigma_\alpha^{+}$ le côté de $E_\alpha$ qui
ne contient pas $V$.

\pagerange{23}{24}
\subsection*{L'arbre des morceaux (pages 23 et 24)}

Soit $S$ l'ensemble des morceaux $U_s$ de $\Pi_I^0$. Disons que le morceau
$U_s$ et la coupure $E_i$ sont \emph{incidents} si $\overline{U_s} \cap E_i
\neq \emptyset$.

\textbf{Lemme (page 23).} \emph{Chaque coupure $E_i$ est incidente à
exactement deux morceaux, l'un dans $\Sigma_i^{+}$, l'autre dans
$\Sigma_i^{-}$. Le graphe $\Gamma_I$ qui a pour sommets les morceaux et pour
arêtes les coupures, l'arête $i$ joignant ses deux morceaux incidents, est
sans arête multiple, et c'est un arbre.}

La page ne démontre que le début\footnote{Elle écrit que $(S, I, R)$ est un
« graphe strict », ce qui signifie a) que chaque $E_i$ touche exactement
deux morceaux, et b) que « $i$ est caractérisé par l'une de ces deux
$U_s$ » ; on lit b) comme l'absence d'arêtes multiples : $i$ est déterminé
par ses deux extrémités. La démonstration qui suit est de nous.}. Voici un
argument complet, qui suit la récurrence des pages 16 à 19. Pour $I = J
\sqcup \{i\}$, le graphe $\Gamma_I$ s'obtient à partir de $\Gamma_J$ en
dédoublant le sommet $U_{J,i}$ en deux sommets $U_{J,i} \cap \Sigma_i^{\pm}$
joints par l'arête $i$, les autres arêtes issues de $U_{J,i}$ allant à l'un
ou à l'autre selon le côté de $E_i$ où se trouve la coupure. En effet :
\begin{itemize}
  \item si $E_j \subset \Sigma_i^{-}$ touche $U_{J,i}$, il touche $U_{J,i}
  \cap \Sigma_i^{-}$, puisque $\Sigma_i^{-}$ est un voisinage ouvert de $E_j$,
  et ne touche pas $U_{J,i} \cap \Sigma_i^{+}$, dont l'adhérence est dans
  $\Sigma_i^{+} \cup E_i$ ;
  \item un morceau $U_t \neq U_{J,i}$ ne touche pas $E_i$, car il est dans
  un côté d'une coupure $E_j$ qui ne contient pas $E_i$, dont l'adhérence est
  dans ce côté réuni à $E_j$ ;
  \item les deux moitiés touchent $E_i$, puisque $U_{J,i}$ est un voisinage
  de $E_i$ et que les deux côtés de $E_i$ l'ont dans leur adhérence.
\end{itemize}
Dédoubler un sommet d'un arbre en une arête, en répartissant les autres
arêtes, donne un arbre ; on part de l'arbre à un sommet ($I = \emptyset$).

La page 24 commence à décrire le morceau incident à $E_i$ du côté
$\Sigma_i^{-}$. Il s'écrit $U = X \smallsetminus \bigcup_{\alpha \in K}
\widehat{\Sigma}_\alpha^{+}$ avec $i \in K$, et $\Sigma_i^{+}$ y est le côté
de $E_i$ qui ne contient pas $U$. Il s'agit de voir que le reste de
l'écriture est déterminé par $i$ et le choix du côté. C'est vrai, et
\[
K \smallsetminus \{i\} = \{ \alpha \neq i \mid E_\alpha \subset
\Sigma_i^{-} \text{ et aucune coupure ne sépare } E_\alpha \text{ de } E_i
\},
\]
$\Sigma_\alpha^{+}$ étant, pour ces $\alpha$, le côté de $E_\alpha$ qui ne
contient pas $E_i$ : ce sont les arêtes de l'arbre qui partent du même
sommet que $i$\footnote{La page s'arrête sur « $K \smallsetminus \{i\} = \{
\alpha \in I \mid E_\alpha \subset \Sigma_\alpha^{-} \}$ », le reste restant
blanc. Telle qu'écrite la condition ne dépend pas de $i$ ; lue avec
$\Sigma_i^{-}$, comme le transcripteur le suggère, elle est fausse dès trois
cordes parallèles $E_1, E_2, E_3$ d'un disque, dans cet ordre : du côté de
$E_1$ qui contient les deux autres, le morceau incident à $E_1$ ne touche que
$E_2$. La condition de non-séparation est de nous.}.

\pagerange{25}{29}
\section*{IV. Dispositions binaires et arbres (pages 25 à 29)}

\pagerange{25}{26}
\subsection*{Dispositions binaires et paires non parallèles (pages 25 et 26)}

Une \emph{disposition binaire} d'un ensemble $X$ est une paire de parties
$\{\Sigma^{+}, \Sigma^{-}\}$, les \emph{hémisphères}, de réunion $X$, dont
aucune ne contient l'autre ; leur intersection $E$, éventuellement vide, est
l'\emph{équateur}. Pour un espace topologique, on demande les hémisphères
fermés, et le plus souvent connexes avec $E \neq \emptyset$ ; alors $X$ est
connexe. Si $X$ est connexe, l'hémisphère ouvert $\Sigma^{\pm,0} = \Sigma^{\pm}
\smallsetminus E$ a $E$ dans son adhérence : sinon il serait ouvert et
fermé, non vide et distinct de $X$\footnote{La première phrase topologique
est d'une lecture incertaine ; on a gardé ce qui s'y démontre.}.

Deux dispositions binaires $\{\Sigma_1^{\pm}\}$, $\{\Sigma_2^{\pm}\}$ sont
\emph{non parallèles} si elles sont distinctes et si l'on peut choisir un
hémisphère de chacune, noté $\Sigma_1^{+}$ et $\Sigma_2^{+}$, avec
$\Sigma_1^{+} \cap \Sigma_2^{+} = \emptyset$\footnote{« Si elles sont
distinctes » est ajouté en interligne ; la démonstration de la page 26 en a
besoin.}. Alors $E_1 \cap E_2 = \emptyset$, $\Sigma_1^{+} \subset
\Sigma_2^{-}$ et $\Sigma_2^{+} \subset \Sigma_1^{-}$, et $\Sigma_{12} =
\Sigma_1^{-} \cap \Sigma_2^{-}$ contient $E_1 \cup E_2$ et n'est pas vide,
même si les deux équateurs le sont : sinon les deux dispositions seraient
deux partitions de $X$ égales à $\{\Sigma_1^{+}, \Sigma_2^{+}\}$. Il en
résulte que le choix de $\Sigma_1^{+}$ et $\Sigma_2^{+}$ est unique. La
chaîne d'inclusions
\[
\Sigma_1^{+} \supset E_1 \hookrightarrow \Sigma_{12} \hookleftarrow E_2
\subset \Sigma_2^{+}
\]
se dessine comme un segment à trois points marqués, et la marge note
$\Sigma_1^{-} = \Sigma_{12} \amalg_{E_2} \Sigma_2^{+}$ et $\Sigma_2^{-} =
\Sigma_{12} \amalg_{E_1} \Sigma_1^{+}$, sommes amalgamées.

Une longue note oblique de la marge signale que les quatre hémisphères ne
sont pas nécessairement distincts --- on peut avoir $\Sigma_1^{+} =
\Sigma_2^{-}$, d'où $E_1 = \Sigma_{12}$ et $E_2 = \emptyset$ --- et ajoute
« strict.\ » devant « non parallèles » : c'est la notion que définit la
page 36\footnote{La note est d'une lecture très incertaine ; on n'en garde
que ce que la page 35 établit.}.

\pagerange{27}{29}
\subsection*{D'un arbre à une famille de décompositions (pages 27 à 29)}

Voici « une façon très générale » d'obtenir des familles de décompositions
binaires deux à deux non parallèles. Soit $\Gamma = (S, A)$ un arbre
\emph{fini}, d'ensemble de sommets $S$ et d'ensemble d'arêtes $A$, et soient
$(U_s)_{s \in S}$, $(E_a)_{a \in A}$ deux familles d'ensembles ; posons
\[
X = \coprod_{s \in S} U_s \ \amalg\ \coprod_{a \in A} E_a .
\]
Pour $a_0 \in A$, l'arbre privé de l'arête $a_0$ a deux composantes
$\Gamma'$, $\Gamma''$ ; soient $S'$, $S''$ leurs sommets, et $A'$, $A''$
leurs arêtes \emph{auxquelles on ajoute $a_0$}. La décomposition binaire
$\Pi_{a_0}$ est $X = X' \cup X''$, avec
\[
X' = \coprod_{s \in S'} U_s \amalg \coprod_{a \in A'} E_a, \qquad
X'' = \coprod_{s \in S''} U_s \amalg \coprod_{a \in A''} E_a, \qquad
X' \cap X'' = E_{a_0} .
\]
La page ne dit pas si $A'$ et $A''$ contiennent $a_0$ ; l'égalité $X' \cap X''
= E_{a_0}$ l'exige\footnote{La finitude de l'arbre est de nous : sans elle,
une composante peut n'avoir aucune feuille, et la condition ($*$) ne suffit
plus.}.
Elle est \emph{propre} ($X' \neq E_{a_0}$ et $X'' \neq E_{a_0}$) pour tout
$a_0$ si et seulement si
\begin{enumerate}
  \item[($*$)] pour toute feuille $s$ de $\Gamma$, $U_s \neq \emptyset$.
\end{enumerate}
En effet, chaque composante $\Gamma'$ contient une feuille de $\Gamma$, et
réciproquement, si $U_s = \emptyset$ pour une feuille $s$, la coupure à son
arête isole $X' = E_{a_0}$.

Pour $a \neq b$, l'hémisphère de $\Pi_a$ qui ne contient pas $E_b$ et celui
de $\Pi_b$ qui ne contient pas $E_a$ sont disjoints ; $\Pi_a$ et $\Pi_b$ sont donc non
parallèles dès qu'elles sont distinctes, et elles le sont si et seulement
si
\begin{enumerate}
  \item[($**$)] pour tout sommet $s$ d'ordre $2$, d'arêtes $a$ et $b$, l'un
  des trois ensembles $U_s$, $E_a$, $E_b$ est non vide\footnote{La condition
  est écrite en partie dans l'interligne et en partie dans la marge, et sa
  lecture est incertaine. On a vérifié qu'ainsi lue elle est exactement ce
  qu'il faut : $\Pi_a = \Pi_b$ exige $E_a = E_b = \emptyset$ et tout ce qui
  est entre $a$ et $b$ vide, ce qui, sous ($*$), force les sommets
  intermédiaires à être d'ordre $2$ et viole ($**$) au premier d'entre
  eux.}.
\end{enumerate}
Ils sont \emph{strictement} non parallèles au sens de la page 36 si et
seulement si l'on a de plus
\begin{enumerate}
  \item[($*{*}*$)] pour tout sommet $s$ d'ordre $2$ et toute arête $a$ qui y
  aboutit, $U_s \neq \emptyset$ ou $E_a \neq \emptyset$,
\end{enumerate}
condition qui entraîne ($**$). La page 29 l'affirme, en ajoutant que le cas
où certains ensembles sont vides « est à examiner un peu de plus
près »\footnote{Le passage est d'une lecture incertaine. On a vérifié
l'énoncé sur les cas de la page 35 : l'égalité $E_a = \Sigma_{12}$ demande
que tout ce qui sépare $a$ de $b$, et $E_b$, soient vides, ce que ($*{*}*$)
interdit au sommet voisin de $b$.}.

\textbf{La réalisation topologique.} La page 28 visualise la construction.
Réalisons $\Gamma$ comme graphe topologique $|\Gamma|$ et marquons le milieu
$a'$ de chaque arête : c'est la subdivision barycentrique $\Gamma^{*}$, de
sommets $S^{*} = S \amalg A'$. Chaque point $a'$ coupe $|\Gamma|$ en deux
composantes, dont les adhérences forment une décomposition binaire
d'équateur $\{a'\}$, et deux telles décompositions sont strictement non
parallèles. Leurs traces sur $S^{*}$ sont des décompositions binaires de
$S^{*}$ ; et se donner une application surjective $X \to S^{*}$ revient à se
donner les deux familles $(U_s)$, $(E_a)$ d'ensembles non vides. La
construction de la page 27 est donc l'image réciproque, par une telle
application, des décompositions de l'arbre lui-même\footnote{La page 28
annonce, entre crochets, « une description complète des propriétés
particulières des $\Gamma^{*}$ » --- chaque arête a une extrémité dans $A'$
et l'autre dans $S$ --- et s'interrompt.}.

C'est la construction inverse de celle de la Version I, qui part des
coupures pour trouver l'arbre ; les deux ensemble disent qu'une famille
finie de découpages deux à deux non parallèles et un arbre sont la même
chose. Dans le cas d'un ensemble fini sans équateurs, c'est ce qu'on appelle
aujourd'hui le \emph{théorème d'équivalence des scissions} de
Buneman\footnote{P.~Buneman, 1971 : une famille de scissions (bipartitions)
d'un ensemble fini deux à deux compatibles --- l'un des quatre croisements
vide, exactement la condition de non-parallélité --- est l'ensemble des
scissions définies par les arêtes d'un arbre. Ce travail a paru dans un
recueil d'archéologie quantitative ; les pages ne le citent pas. Le rapprochement, et l'affirmation que les deux constructions
sont inverses, sont de nous.}.

\pagerange{31}{39}
\section*{V. Décompositions $n$-aires : la Version II (pages 31 à 39)}

La Version II, titrée et numérotée, reprend tout avec des
hémisphères fermés et sans supposer $X$ connexe.

\pagerange{31}{32}
\subsection*{1.1. Décompositions binaires (pages 31 et 32)}

Une \emph{décomposition binaire} d'un espace topologique $X$ est une paire
$\Pi = (\Sigma^{\varepsilon})_{\varepsilon \in \omega}$, indexée par un
ensemble $\omega$ à deux éléments, de parties \emph{fermées} de $X$ telles
que
\[
X = \bigcup_{\varepsilon \in \omega} \Sigma^{\varepsilon} \quad (1), \qquad
\Sigma^{\varepsilon} \neq \emptyset \quad (2), \qquad
\Sigma^{\varepsilon} \neq \Sigma^{\varepsilon'} \text{ pour } \varepsilon
\neq \varepsilon' \quad (3).
\]
Il en résulte $\operatorname{card} X \geqslant 2$. On n'exclut pas que
$\Sigma^{\varepsilon'} \subset \Sigma^{\varepsilon}$, ce qui équivaut par
(1) à $\Sigma^{\varepsilon} = X$. Les $\Sigma^{\varepsilon}$ sont les
\emph{hémisphères}, et
\[
E = \bigcap_{\varepsilon \in \omega} \Sigma^{\varepsilon} \qquad (4)
\]
l'\emph{équateur}. Si $E = \emptyset$, les hémisphères sont complémentaires,
donc ouverts et fermés, et $X$ n'est pas connexe. D'où (5) : si $X$ est
connexe, $E \neq \emptyset$ ; et si les hémisphères sont connexes, $X$ est
connexe si et seulement si $E \neq \emptyset$.

Pour $\omega = \{\varepsilon, \varepsilon'\}$, on pose
\[
\Sigma^{0\varepsilon} = \Sigma^{\varepsilon} \smallsetminus E = X
\smallsetminus \Sigma^{\varepsilon'} \qquad (6),
\]
ouvert, d'où la partition $X = \Sigma^{0\varepsilon} \sqcup E \sqcup
\Sigma^{0\varepsilon'}$ (7), et
\[
\overline{\Sigma^{0\varepsilon}} \cap E = \overline{\Sigma^{0\varepsilon}}
\cap \Sigma^{\varepsilon'} = \overline{\Sigma^{0\varepsilon}} \smallsetminus
\Sigma^{0\varepsilon} = \dot{\Sigma}^{0\varepsilon} \qquad (8),
\]
la frontière de l'hémisphère ouvert\footnote{La page commence ensuite une
condition sur ces frontières (« on s'intéresse à la cond. »), qu'elle barre.
C'est la condition qui, dans la Version I, vient de la connexité.}.

\pagerange{32}{36}
\subsection*{1.2. Paires non parallèles (pages 32 à 36)}

Deux décompositions binaires $\Pi_1 = \{\Sigma_1^{+}, \Sigma_1^{-}\}$ et
$\Pi_2 = \{\Sigma_2^{+}, \Sigma_2^{-}\}$ (9) sont \emph{non parallèles} si
elles sont distinctes et si l'on peut nommer leurs hémisphères de sorte que
\[
\Sigma_1^{+} \cap \Sigma_2^{+} = \emptyset \qquad (10).
\]
Cette condition équivaut à $\Sigma_1^{+} \subset \Sigma_2^{0-}$, et à
$\Sigma_2^{+} \subset \Sigma_1^{0-}$ ; elle entraîne les inclusions
\[
\Sigma_1^{+} \subset \Sigma_2^{-} \quad (11), \qquad \Sigma_2^{+} \subset
\Sigma_1^{-} \quad (11 \text{ bis}),
\]
mais n'équivaut à aucune d'elles\footnote{La page écrit que (10)
« équivaut » à (11), « ou encore » à (11 bis). Avec des hémisphères fermés,
(11) n'exclut pas que $\Sigma_1^{+}$ touche l'équateur $E_2$. Exemple (de
nous) : dans un disque, $\Sigma_2^{-}$ un demi-disque fermé, $\Sigma_1^{+}$ un
petit disque fermé contenu dans $\Sigma_2^{-}$ et tangent au diamètre $E_2$,
$\Sigma_1^{-}$ l'adhérence de son complémentaire ; (11) est vraie, et aucune
paire d'hémisphères n'est disjointe. Dans la Version I, où $\Sigma$ désigne
les composantes ouvertes, l'équivalence de la page 14 est juste.}. La marge
ajoute que (10) donne $E_1 \subset \Sigma_2^{0-}$, $E_2 \subset
\Sigma_1^{0-}$, et $E_1 \cap E_2 = \emptyset$.

Posons $\Sigma_{12} = \Sigma_1^{-} \cap \Sigma_2^{-}$, $E_1 = \Sigma_1^{+}
\cap \Sigma_1^{-}$, $E_2 = \Sigma_2^{+} \cap \Sigma_2^{-}$ (13). On a le
diagramme d'inclusions de parties fermées (14) :

\begin{tikzcd}[column sep=small, row sep=small]
 & & X & & \\
 & \Sigma_{2}^{-} \arrow[ur] & & \Sigma_{1}^{-} \arrow[ul] & \\
 \Sigma_{1}^{+} \arrow[ur] & & \Sigma_{12} \arrow[ul] \arrow[ur] & &
 \Sigma_{2}^{+} \arrow[ul] \\
 & E_{1} \arrow[ul] \arrow[ur] & & E_{2} \arrow[ul] \arrow[ur] &
\end{tikzcd}

où les trois carrés sont \emph{cartésiens et cocartésiens} : dans chacun, le
sommet du bas est l'intersection et le sommet du haut la réunion des deux
autres,
\[
\begin{array}{ll}
E_1 = \Sigma_1^{+} \cap \Sigma_{12}, & \Sigma_2^{-} = \Sigma_1^{+} \cup \Sigma_{12}, \\
E_2 = \Sigma_{12} \cap \Sigma_2^{+}, & \Sigma_1^{-} = \Sigma_{12} \cup \Sigma_2^{+}, \\
\Sigma_{12} = \Sigma_1^{-} \cap \Sigma_2^{-}, & X = \Sigma_2^{-} \cup \Sigma_1^{-}
\end{array}
\qquad (15).
\]
Ce sont des vérifications directes à partir de (11) et (11 bis)\footnote{Par
exemple $\Sigma_1^{+} \cup \Sigma_{12} = (\Sigma_1^{+} \cup \Sigma_1^{-})
\cap (\Sigma_1^{+} \cup \Sigma_2^{-}) = \Sigma_2^{-}$. Une réunion de deux
fermés recollés le long de leur intersection est aussi une somme amalgamée
d'espaces topologiques : les carrés sont cocartésiens dans les espaces, pas
seulement dans l'ensemble ordonné des parties.}. De plus (16) :
$\Sigma_1^{+} \cap \Sigma_2^{-} = \Sigma_1^{+} \neq \emptyset$,
$\Sigma_2^{+} \cap \Sigma_1^{-} = \Sigma_2^{+} \neq \emptyset$, et
$\Sigma_{12} \neq \emptyset$ ; pour la dernière, si $\Sigma_{12}$ était vide,
les équateurs qu'il contient le seraient, $\Pi_1$ et $\Pi_2$ seraient des
partitions, et $X = \Sigma_1^{+} \amalg \Sigma_2^{+}$ donnerait $\Pi_1 = \Pi_2$.

\textbf{Proposition (page 34).} \emph{Si $\Pi_1$ et $\Pi_2$ sont non
parallèles, il existe un unique couple $(\Sigma_1^{+}, \Sigma_2^{+}) \in
\Pi_1 \times \Pi_2$ d'hémisphères disjoints.} Les trois autres couples se
rencontrent, par (16).

Soit $C = \Pi_1 \amalg \Pi_2 = \omega_1 \amalg \omega_2$ (17), ensemble à
quatre éléments, et $C \to \mathfrak{P}_{\text{fermé}}(X)$ (18) l'application
évidente, d'image $\{\Sigma_1^{+}, \Sigma_1^{-}, \Sigma_2^{+}, \Sigma_2^{-}\}$
(19). Elle est injective si et seulement si les quatre hémisphères sont
distincts. On sait déjà ($*$) $\Sigma_i^{+} \neq \Sigma_i^{-}$, par (3), et
($**$) $\Sigma_1^{+} \neq \Sigma_2^{+}$, parties disjointes non vides. On lit
sur (14) :
\[
\begin{array}{ll}
\text{a) } (\Sigma_1^{+} = \Sigma_2^{-}) \Longleftrightarrow E_1 = \Sigma_{12} & (\Rightarrow E_2 = \emptyset), \\
\text{b) } (\Sigma_2^{+} = \Sigma_1^{-}) \Longleftrightarrow E_2 = \Sigma_{12} & (\Rightarrow E_1 = \emptyset)
\end{array}
\qquad (*{*}*),
\]
et les deux ne peuvent avoir lieu ensemble, qui donneraient $\Sigma_{12} =
E_1 = E_2 = \emptyset$, contre (16)\footnote{Dans la première relation, le
signe entre $\Sigma_1^{+}$ et $\Sigma_2^{-}$ est un $=$ repassé en $\neq$, ou
l'inverse ; la seconde et la suite demandent $=$.}. Enfin
\[
(\Sigma_1^{-} = \Sigma_2^{-}) \Longleftrightarrow (\Sigma_1^{-} =
\Sigma_2^{-} = X) \Longleftrightarrow \Sigma_{12} = X \qquad (*{*}{*}*),
\]
cas qui exclut chacun des deux précédents, puisque $E_1 = \Sigma_{12} = X$
contredirait (3).

\textbf{Proposition (page 35).} \emph{Les quatre hémisphères sont distincts
si et seulement si l'on n'est dans aucun des trois cas suivants, qui
s'excluent mutuellement :}
\[
\text{a) } E_1 = \Sigma_{12}; \qquad \text{a') } E_2 = \Sigma_{12}; \qquad
\text{b) } \Sigma_{12} = X, \text{ i.e. } \Sigma_1^{-} = \Sigma_2^{-} = X
\qquad (20).
\]
\emph{Dans chacun, l'ensemble (19) a trois éléments.} Dans le cas b), on a
$E_1 = \Sigma_1^{+}$ et $E_2 = \Sigma_2^{+}$.

\textbf{Corollaire (page 36).} \emph{Si les quatre hémisphères sont
distincts et si $\Sigma_1^{-} \neq X$, $\Sigma_2^{-} \neq X$, les seules
inclusions entre eux, hors les égalités, sont $\Sigma_1^{+} \subset
\Sigma_2^{-}$ et $\Sigma_2^{+} \subset \Sigma_1^{-}$.} Les inclusions
inverses donneraient des égalités ; $\Sigma_1^{+} \subset \Sigma_1^{-}$
donnerait $\Sigma_1^{-} = X$, et $\Sigma_1^{-} \subset \Sigma_1^{+}$
donnerait $\Sigma_1^{+} = X$, impossible puisque $\Sigma_2^{+}$ est non vide
et disjoint de $\Sigma_1^{+}$\footnote{La page conclut « on aurait
$\Sigma_1^{-} = X$ » pour les deux inclusions ; pour la seconde c'est
$\Sigma_1^{+} = X$, exclu pour la raison dite.} ; de même pour l'indice 2 ;
$\Sigma_1^{+}$ et $\Sigma_2^{+}$ sont disjoints et non vides ; et
$\Sigma_1^{-} \subset \Sigma_2^{-}$ donnerait $\Sigma_2^{-} \supset
\Sigma_1^{-} \cup \Sigma_1^{+} = X$. La marge reformule : ce sont les
seules inclusions non triviales si et seulement si les quatre sont distincts
et $\Sigma_1^{-} \neq X$, $\Sigma_2^{-} \neq X$.

\textbf{Définition (page 36).} $\Pi_1$ et $\Pi_2$ sont \emph{strictement non
parallèles} si elles sont non parallèles, si les quatre hémisphères sont
distincts, et si $\Sigma_1^{-} \neq X$, $\Sigma_2^{-} \neq X$ ; autrement
dit, si les seules inclusions entre les quatre sont $\Sigma_1^{+} \subset
\Sigma_2^{-}$ et $\Sigma_2^{+} \subset \Sigma_1^{-}$ ; autrement dit encore,
\[
\Sigma_1^{-} \neq X, \quad \Sigma_2^{-} \neq X \quad (\Pi_1, \Pi_2
\text{ \emph{propres}}), \qquad E_1 \neq \Sigma_{12}, \quad E_2 \neq
\Sigma_{12} \qquad (21).
\]

\pagerange{37}{38}
\subsection*{1.3. Familles non parallèles et leur ensemble ordonné (pages 37 et 38)}

Que la condition stricte soit remplie ou non, on munit l'ensemble à quatre
éléments $C = \{\varepsilon_1^{+}, \varepsilon_1^{-}, \varepsilon_2^{+},
\varepsilon_2^{-}\}$ (17) de l'ordre dont les seules relations strictes sont
\[
\varepsilon_1^{+} < \varepsilon_2^{-}, \qquad \varepsilon_2^{+} <
\varepsilon_1^{-} \qquad (22),
\]
que la page dessine :

\begin{tikzcd}[column sep=small, row sep=small]
\varepsilon_{1}^{-} & \varepsilon_{2}^{-} \\
\varepsilon_{1}^{+} \arrow[ur] & \varepsilon_{2}^{+} \arrow[ul]
\end{tikzcd}

Soit maintenant $(\Pi_\alpha)_{\alpha \in A}$, $\Pi_\alpha =
\{\Sigma_\alpha^{\varepsilon}\}_{\varepsilon \in \omega_\alpha}$ (23), un
ensemble de décompositions binaires de $X$ deux à deux non parallèles (24).
\textbf{Définition (page 37).} On dit que la famille est \emph{non
parallèle}, et, si $\operatorname{card} A = n$, que c'est une
\emph{décomposition $n$-aire} non parallèle de $X$ ;
terminologie analogue pour « strictement »\footnote{La marge note que c'est
un élément de $\mathfrak{P}(\mathfrak{P}_2(\mathfrak{P}(X)))$ : un ensemble
de paires de parties.}. Considérons
\[
C = \coprod_{\alpha \in A} \omega_\alpha \qquad (25),
\]
l'involution sans point fixe $\sigma$ qui échange les deux éléments de chaque
$\omega_\alpha$, et la relation définie, pour $x \in \omega_\alpha$ et $y \in
\omega_\beta$, par
\[
x < y \ \Longleftrightarrow\ \alpha \neq \beta \ \text{ et } \
\Sigma_\alpha^{x} \cap \Sigma_\beta^{\sigma y} = \emptyset \qquad (26),
\]
c'est-à-dire par l'ordre (22) de la paire $\{\Pi_\alpha, \Pi_\beta\}$.

\textbf{C'est un ordre}, ce que la marge demande de vérifier\footnote{« Il
faut vérifier que c'est transitif … d'ordre », en marge. La vérification
est de nous.}. L'antisymétrie vient de l'unicité du couple disjoint
(proposition de la page 34) : $x < y$ et $y < x$ désigneraient comme couple
disjoint de $\Pi_\alpha \times \Pi_\beta$ à la fois $(x, \sigma y)$ et
$(\sigma x, y)$. Pour la transitivité, $x < y < z$ donne $\Sigma^{x} \subset
X \smallsetminus \Sigma^{\sigma y} \subset \Sigma^{y}$ et $\Sigma^{y} \cap
\Sigma^{\sigma z} = \emptyset$, donc $\Sigma^{x} \cap \Sigma^{\sigma z} =
\emptyset$ ; et $x$, $z$ ne sont pas dans le même $\omega_\alpha$, car $z = x$
contredirait l'antisymétrie et $z = \sigma x$ donnerait $\Sigma^{x} =
\emptyset$.

L'involution $\sigma$ \emph{renverse l'ordre} : $x < y$ entraîne $\sigma y <
\sigma x$, la condition (26) étant symétrique. L'application
\[
\varphi \colon C \to \mathfrak{P}_{\text{fermé}}(X), \qquad (\alpha, x)
\mapsto \Sigma_\alpha^{x} \qquad (27)
\]
est croissante, et c'est un isomorphisme d'ensembles ordonnés sur son image
(ordonnée par inclusion) si et seulement si la famille est strictement non
parallèle. Cela résulte du corollaire de la page 36 appliqué à chaque paire,
et, pour $x$ et $\sigma x$, qui sont incomparables dans $C$, de ce que
$\Sigma^{x}$ et $\Sigma^{\sigma x}$ ne sont comparables que si l'un vaut $X$,
ce que « propre » exclut\footnote{La page affirme l'énoncé sans
démonstration ; celle-ci est de nous.}.

\pagerange{38}{39}
\subsection*{Les quartettes ordonnés (pages 38 et 39)}

« Inversement », la page pose la définition abstraite de ce qu'elle vient
d'obtenir.

\textbf{Définition (page 38).} Un \emph{quartette ordonné} est un ensemble
fini $C$ muni d'une involution sans point fixe $\sigma$ et d'une relation
d'ordre, tels que :
\begin{enumerate}
  \item[a)] pour $x, y \in C$, les éléments $x$ et $y$ sont comparables si et
  seulement si $\sigma x$ et $y$ ne le sont pas ;
  \item[b)] $\sigma$ renverse l'ordre : $x < y$ entraîne $\sigma y < \sigma
  x$.
\end{enumerate}
L'ensemble ordonné $C$ d'une famille non parallèle est un quartette ordonné :
pour $\alpha \neq \beta$, des quatre paires $\{x, y\}$, $\{\sigma x, y\}$,
$\{x, \sigma y\}$, $\{\sigma x, \sigma y\}$, exactement deux, échangées par
$\sigma$, sont comparables ; et $x$ est comparable à lui-même, non à $\sigma
x$\footnote{« Fini » est une lecture incertaine, et « ordonné » est biffé
puis, semble-t-il, récrit. La condition b) est biffée sur la page, mais son
contenu revient dans le NB qui suit (« $\sigma$ est un antiautomorphisme de
l'ens.\ ordonné $C$ »), sans qu'on puisse lire s'il est posé ou déduit. Il
faut le poser : a) seul ne l'entraîne pas. Contre-exemple (de nous) : $C =
\{x, \sigma x, y, \sigma y\}$ avec pour seules relations $x < y$ et $\sigma x
< \sigma y$ vérifie a), et $\sigma$ n'y renverse pas l'ordre.}.

Ce sont exactement, dans le vocabulaire d'aujourd'hui, les « pocsets » finis
dont deux murs quelconques sont emboîtés ; et un tel pocset est l'ensemble
des demi-arbres d'un arbre fini, ordonné par inclusion\footnote{Le terme
« pocset » (ensemble ordonné muni d'une involution qui renverse l'ordre,
$x$ et $\sigma x$ incomparables) est de M.~Roller (1998) ; la construction
d'un arbre à partir d'une famille emboîtée de parties est de M.~J.~Dunwoody
(1979), et l'on retrouve le cube complexe de M.~Sageev (1995) quand on lève
l'emboîtement. Les pages ne citent aucun de ces travaux. L'identification des quartettes ordonnés aux
pocsets emboîtés, et donc aux arbres finis, est de nous. Elle donnerait la
réciproque qu'annonce « Inversement » : tout quartette ordonné provient d'une
famille strictement non parallèle, celle de la réalisation topologique de
son arbre (page 28). La page ne l'énonce pas.}. La réciproque de la
construction des pages 27 à 29 est ainsi à portée : c'est là, semble-t-il,
que la rédaction voulait aller.

\textbf{Proposition (page 39), telle qu'énoncée.} \emph{Soit $(C, \sigma, <)$
un quartette ordonné. Pour tout $x \in C$, les ensembles $C_{\geqslant x} =
\{y \mid y \geqslant x\}$ et $C_{\leqslant x} = \{y \mid y \leqslant x\}$
sont totalement ordonnés.}

\emph{L'énoncé est faux.} Prenons trois cordes disjointes d'un disque qui
entourent un triangle, et les trois décompositions binaires qu'elles
définissent, de calottes $c_1$, $c_2$, $c_3$ et de grands côtés $c_1'$,
$c_2'$, $c_3'$. La famille est strictement non parallèle, et son quartette
est $C = \{c_1, c_2, c_3, c_1', c_2', c_3'\}$, avec $\sigma c_i = c_i'$ et
pour seules relations $c_i < c_j'$ ($i \neq j$). Alors $C_{\geqslant c_1} =
\{c_1, c_2', c_3'\}$, et $c_2'$, $c_3'$ sont incomparables. Ce qui est vrai
est plus faible :
\begin{itemize}
  \item \emph{tout intervalle $[x, z] = \{y \mid x \leqslant y \leqslant z\}$
  est totalement ordonné} ;
  \item \emph{$C_{\geqslant x}$ est totalement ordonné pour tout $x$ si et
  seulement si l'arbre du quartette est une chaîne}, c'est-à-dire si toutes
  les coupures sont « en file », comme des cordes parallèles au sens usuel\footnote{Le
  contre-exemple et les deux énoncés vrais sont de nous. Le graphe du
  quartette de trois cordes autour d'un triangle est un tripode, qu'on
  trouve dessiné, en « Y » épais, page 11.}.
\end{itemize}
Pour le premier point : si $x \leqslant y, y' \leqslant z$ avec $y$, $y'$
incomparables, a) rend $\sigma y$ et $y'$ comparables ; $\sigma y \leqslant
y' \leqslant z$ rendrait $y$ et $\sigma y$ tous deux comparables à $z$, et
$y' \leqslant \sigma y$ rendrait $y$ et $\sigma y$ tous deux comparables à
$x$, ce que a) interdit dans les deux cas.

La démonstration de la page suit exactement ce chemin, et s'arrête là où il
n'y a plus rien à trouver. Elle se ramène, par dualité, à voir que $x < y$,
$x < z$, $y \neq z$ entraînent $y$, $z$ comparables. Elle note que $y$ et $z$
ne sont pas conjugués, puis suppose $y$, $z$ incomparables : alors $\sigma y$
et $z$ sont comparables ; $z \leqslant \sigma y$ est exclu, car il donnerait
$x \leqslant \sigma y$ ; donc $\sigma y < z$, d'où $\sigma z < y$ --- et la
phrase s'arrête sur « d'où $\sigma z <$ »\footnote{La page invoque deux fois
une condition « (28 a) », qu'on ne trouve pas sur la page 38 transcrite. Les
deux usages sont des conséquences de a) et de b) : $x < y$ et $x \leqslant
\sigma y$ rendraient $\sigma x$ comparable à $y$. On identifie donc (28 a) à
la condition a) ; le numéro n'est pas lisible sur la page.}. Dans l'exemple
des trois cordes, avec $x = c_1$, $y = c_2'$, $z = c_3'$, on a bien $\sigma y
= c_2 < c_3' = z$ et $\sigma z = c_3 < c_2' = y$, sans contradiction.

\pagerange{40}{52}
\section*{VI. Relations d'équivalence non parallèles (pages 40 à 52)}

La chemise de la page 40 porte « relations d'équi.\ non parallèles ». Les
pages qui suivent laissent l'espace topologique et ne gardent que les
extrémités des coupures : des points rangés sur un cercle.

\pagerange{41}{42}
\subsection*{Polygones et segments (pages 41 et 42)}

Soit $S$ un ensemble fini de cardinal $N \geqslant 3$, muni d'une
\emph{structure polygonale} ou \emph{bi-circulation} $\{u, u^{-1}\}$, où $u$
est une permutation circulaire de $S$\footnote{La page ne dit pas que $u$
est circulaire, c'est-à-dire transitif ; c'est ce que veulent dire
« polygonale » et tout ce qui suit. C'est la troisième description d'un
polygone combinatoire du dossier 89 (pages 19 et 20), où une orientation
est le choix de l'une des deux permutations.}. Une partie $T$ de cardinal $i$
est un \emph{segment} si $0 < i < N$ et s'il existe $s_0 \in S$ avec
\[
T = \{s_0, u s_0, \ldots, u^{i-1} s_0\} ;
\]
avec $t_0 = u^{i-1} s_0$ on a aussi $T = \{t_0, v t_0, \ldots, v^{i-1} t_0\}$,
où $v = u^{-1}$. \emph{Exemples} : les singletons, les paires de sommets
voisins (les « parties-arêtes »). Propriétés :
\begin{itemize}
  \item $T$ est un segment si et seulement si $S \smallsetminus T$ en est un ;
  \item l'intersection de deux segments contenus dans un même segment est un
  segment si elle n'est pas vide ;
  \item la réunion de deux segments qui se rencontrent et ne recouvrent pas
  $S$ est un segment ;
  \item les \emph{extrémités} de $T$, $\partial T = \{s_0, t_0\}$, sont
  distinctes si et seulement si $T$ n'est pas un singleton ; une orientation
  (le choix de $u$) les distingue en origine et extrémité.
\end{itemize}
Une orientation étant choisie, $(s_0, i) \mapsto \{s_0, \ldots, u^{i-1}
s_0\}$ est une bijection de $S \times [1, N-1]$ sur l'ensemble des
segments, qui en a donc $N(N-1)$\footnote{La page écrit « l'ens.\
$(S,[1,N-1]_{\mathbb{N}})$ » pour le produit. Le décompte est de nous.}.

\textbf{Composantes d'une partie.} Soit $S' \subsetneq S$. Deux segments
contenus dans $S'$ qui se rencontrent ont pour réunion un segment contenu
dans $S'$. Tout segment contenu dans $S'$ est donc contenu dans un plus grand
segment contenu dans $S'$ ; ces segments maximaux, les \emph{segments
composants} de $S'$, forment une partition de $S'$ : ce sont les suites
maximales de sommets consécutifs de $S'$. Si l'on réalise le polygone par un
cercle $X$ portant les points de $S$, les segments composants de $S'$
correspondent aux composantes connexes de $X \smallsetminus (S \smallsetminus
S')$ qui rencontrent $S'$\footnote{La marge introduit
$\operatorname{Segm}_{S'}(S)$, l'ensemble des segments de $S$ contenus dans
$S'$ ; elle est en grande partie illisible.}.

\pagerange{42}{43}
\subsection*{Parties non parallèles (pages 42 et 43)}

\textbf{Proposition (pages 42 et 43).} \emph{Soient $A$, $B$ deux parties
disjointes de $S$, non vides et distinctes de $S$. Les conditions suivantes
sont équivalentes :}
\begin{enumerate}
  \item[a)] \emph{il existe des segments disjoints $\overline{A} \supset A$ et
  $\overline{B} \supset B$ ;}
  \item[b)] \emph{pour $a \neq a'$ dans $A$ et $b \neq b'$ dans $B$, les
  sommets $a$ et $a'$ sont voisins dans le polygone à quatre sommets $\{a, a',
  b, b'\}$ muni de la bi-circulation induite --- ce qui équivaut à ce que $b$
  et $b'$ le soient, ou à ce que $a$ et $a'$ n'y soient pas opposés ;}
  \item[c)] \emph{$B$ est contenu dans un seul segment composant de $S
  \smallsetminus A$ ;}
  \item[c')] \emph{$A$ est contenu dans un seul segment composant de $S
  \smallsetminus B$.}
\end{enumerate}
\emph{On dit alors que $A$ et $B$ sont non parallèles.}

En termes de cordes, b) dit que les cordes $aa'$ et $bb'$ ne se croisent
jamais. Les implications se vérifient ainsi\footnote{La page énonce les
équivalences sans les démontrer ; l'argument est de nous.}. De a) à c) :
$\overline{B}$ est un segment contenu dans $S \smallsetminus A$, donc dans
l'un de ses segments composants. De c) à a) : si $B$ est dans le composant
$K$, on prend $\overline{B} = K$ et $\overline{A} = S \smallsetminus K$. De c)
à b) : si $b$ et $b'$ sont dans le même composant de $S \smallsetminus A$, un
des deux arcs qui les joignent ne contient aucun point de $A$. De b) à c) :
si $b$ et $b'$ sont dans deux composants différents, chacun des deux arcs
qui les joignent contient un point de $A$, et ces deux points alternent avec
$b$, $b'$. La symétrie de a) donne c').

Il y a un plus petit couple de segments $(\overline{A}, \overline{B})$ comme
dans a) : $A$ est dans un composant $K'$ de $S \smallsetminus B$, qui est
totalement ordonné par $u$, et l'on prend pour $\overline{A}$ le plus petit
segment de $K'$ contenant $A$ ; de même pour $\overline{B}$ ; ces deux
segments sont disjoints, car deux arcs d'un cercle dont aucun ne contient
les extrémités de l'autre ne se rencontrent pas s'ils ne recouvrent pas le
cercle\footnote{La page s'arrête sur « il existe un plus petit… », et la
marge parle des « plus petits » segments. La construction est de nous.}.

\emph{Exemple.} Si $A$ ou $B$ est un singleton, la condition b) est vide :
un point est non parallèle à toute partie qui ne le contient pas. En
général, $A$ et $B$ sont non parallèles si et seulement si toute paire
$A' \in \mathfrak{P}_2(A)$ est non parallèle à toute paire $B' \in
\mathfrak{P}_2(B)$\footnote{La page ajoute entre parenthèses « (plus
généralement si $A \cap B$ est un segment) », qu'on ne sait pas
interpréter ; on ne la reprend pas. L'exemple des singletons est à demi
illisible.}.

\textbf{La marge d).} Si les points de $S$ sont placés, dans l'ordre,
sur une courbe convexe du plan --- en position convexe ---, et si $P_A$
désigne l'enveloppe convexe de $A$, alors $A$ et $B$ sont non parallèles si
et seulement si $P_A \cap P_B = \emptyset$. C'est l'image qui a donné leur
nom actuel à ces parties : on les dit \emph{non croisées}\footnote{La note
marginale est écrite en rotation et en partie surchargée ; on n'en garde que
l'énoncé qui ressort. « Réalisé comme partie convexe du bord d'un disque » est
une lecture incertaine.}.

\pagerange{44}{44}
\subsection*{Structure à étudier (page 44)}

\emph{Structure à étudier} : les partitions $(S_\alpha)_{\alpha \in J}$ de
$S$ en parties deux à deux non parallèles, qu'il appelle \emph{partitions non
parallèles}\footnote{La marge porte un exemple, « partitions … (au moins
deux) … de $S$ », en grande partie illisible.}. Ce sont les \emph{partitions non croisées} d'un cycle, étudiées
sous ce nom par G.~Kreweras en 1972\footnote{G.~Kreweras, « Sur les
partitions non croisées d'un cycle », \emph{Discrete Mathematics} 1 (1972).
Elles sont comptées par les nombres de Catalan et forment un treillis ; la
page ne compte ni n'ordonne, et ne cite pas ce travail. Le vocabulaire « non croisé » est invariant par renversement du
cycle, comme la bi-circulation l'exige.}.

\emph{Intuition} : à un ensemble bi-circulé de cardinal $N$ muni d'une
partition non parallèle correspondrait, « essentiellement », un ensemble
bi-circulé de cardinal pair $N' \leqslant N$ muni d'une partition non
parallèle en classes à deux éléments --- une famille de cordes sans
croisement qui en couvre tous les points\footnote{L'intuition n'est pas
précisée, et un mot qui en limiterait la portée est illisible. On connaît
aujourd'hui une bijection entre partitions non croisées d'un $N$-gone et
appariements non croisés d'un $2N$-gone, où chaque point est dédoublé ; elle
ne réalise pas l'inégalité $N' \leqslant N$ de la page, et l'on ne sait pas
quelle correspondance il avait en vue.}. Sous ce texte, un grand cercle, sans
marque.

\pagerange{45}{47}
\subsection*{Familles emboîtées de segments et leurs cœurs (pages 45 à 47)}

Soit $\Sigma$ un ensemble de segments de $S$ tel que :
\begin{enumerate}
  \item[($\sigma_1$)] deux éléments de $\Sigma$ qui se rencontrent sont
  emboîtés ;
  \item[($\sigma_2$)] si $A \subsetneq B$ dans $\Sigma$, alors $A$ ne contient
  aucune extrémité de $B$.
\end{enumerate}
La condition ($\sigma_1$) fait de $\Sigma$ une famille \emph{laminaire}. Pour
$A \in \Sigma$, posons
\[
A^{*} = A \smallsetminus \bigcup_{B \in \Sigma,\ B \subsetneq A} B,
\]
le \emph{cœur} de $A$. Par ($\sigma_2$), $A^{*} \supset \partial A$, donc
$A^{*} \neq \emptyset$.

\textbf{Proposition (marge de la page 45).} \emph{Si $A \neq B$ dans
$\Sigma$, alors $A^{*} \cap B^{*} = \emptyset$ ; en particulier $A \mapsto
A^{*}$ est injectif. De plus $\bigcup_{A \in \Sigma} A = \bigcup_{A \in
\Sigma} A^{*}$.} Si $A \cap B = \emptyset$, c'est clair ; sinon, par
($\sigma_1$), disons $A \subsetneq B$, et $B^{*}$ ne rencontre pas $A$ par
définition. Pour la seconde assertion, un point $x$ d'un élément de $\Sigma$
appartient au cœur du plus petit élément de $\Sigma$ qui le contient, les
éléments qui contiennent $x$ formant une chaîne\footnote{La marge dit « si
$x \in \bigcup A$, alors $\Sigma(x)$ … non vide … $A \in \Sigma(x)$ minimal
… $x \in A^{*}$ », avec des mots illisibles ; on a complété.}.

\textbf{Conséquences (page 46).} Deux cœurs distincts sont non parallèles :
si $A$ et $B$ sont disjoints, $A$ et $B$ sont des segments disjoints qui les
contiennent ; si $A \subsetneq B$, ce sont $A$ et $S \smallsetminus A$. Si les
cœurs recouvrent $S$, ils forment donc une partition non parallèle de
$S$\footnote{La page qualifie cette partition d'un mot lu, avec doute,
« par-dessus ». La non-parallélité des cœurs, qu'elle n'établit pas ici,
est affirmée page 51 (« $A^{*}$ et $B^{*}$ non parallèles ») ; la
démonstration est de nous.}. Sinon, posons
\[
Z = S \smallsetminus \bigcup_{A \in \Sigma} A^{*} = S \smallsetminus
\bigcup_{A \in \Sigma} A, \qquad \widehat{\Sigma} = \Sigma \cup \{Z\},
\qquad Z^{*} = Z ;
\]
$Z$ n'est pas en général un segment, mais il est non parallèle à chaque
cœur, qui est contenu dans un élément maximal $M$ de $\Sigma$ quand $Z$ est
dans le segment $S \smallsetminus M$. La famille $\widehat{\Sigma}^{*} =
\{A^{*} \mid A \in \widehat{\Sigma}\}$ est alors une partition non parallèle
de $S$\footnote{La page note $R$ l'ensemble résiduel, et $\widehat{\Sigma} =
\Sigma$ si $R = \emptyset$ ; on le note $Z$ pour réserver $R$ à la relation
d'équivalence des pages 47 à 50.}.

Si $Z$ est lui-même un segment, $\widehat{\Sigma}$ satisfait encore
($\sigma_1$) et ($\sigma_2$), et $Z$ y est isolé ; la même partition provient
alors de deux familles. Pour l'éviter, on impose
\begin{enumerate}
  \item[($\sigma_3$)] $Z = S \smallsetminus \bigcup_{A \in \Sigma} A$ n'est
  pas un segment
\end{enumerate}
(c'est le cas si $Z = \emptyset$, qui n'est pas un segment).

\textbf{La famille ne se retrouve pas à partir de la partition (page 47).}
Peut-on reconstruire $A \in \Sigma$ à partir de $\widehat{\Sigma}^{*}$ ? Le
segment $A$ contient $A^{*}$ et en est un segment minimal, puisque $A^{*}
\supset \partial A$ ; mais cela ne le caractérise pas --- le segment
complémentaire de l'intérieur de $A$ contient aussi $\partial A$. Il a en
outre la propriété que tout cœur $B^{*}$ est contenu dans $A$ ou disjoint de
$A$. Même ainsi, et même sous ($\sigma_1$) à ($\sigma_3$), la réponse est
non. Soit $S$ l'hexagone $a_1, b, a_2, a_2', b', a_1'$ --- trois cordes
parallèles $a_1 a_1'$, $b b'$, $a_2 a_2'$ tracées dans un cercle ---, et $R$
la partition en $\{a_1, a_1'\}$, $\{b, b'\}$, $\{a_2, a_2'\}$. La famille
$\Sigma_1 = \{\{a_1, a_1'\}, \{a_2, a_2'\}, \{a_1, a_1', b, b'\}\}$ et sa
symétrique $\Sigma_2 = \{\{a_1, a_1'\}, \{a_2, a_2'\}, \{a_2, a_2', b, b'\}\}$
satisfont ($\sigma_1$) à ($\sigma_3$) et ont la même partition $R$ de
cœurs\footnote{On a vérifié l'exemple de la page : les deux paires $\{a_i,
a_i'\}$ sont des segments de l'hexagone, le grand segment de $\Sigma_1$ est
$b', a_1', a_1, b$, et $\{a_1, a_1'\}$ en est intérieur ; son cœur est $\{b,
b'\}$, et $Z = \emptyset$. La suite, un « exemple plus général », est
barrée.}. C'est la même rangée de cordes parallèles, en plus court, que celle
que la page 9 marque d'un « ? ».

\pagerange{48}{50}
\subsection*{Relations non parallèles et anglets (pages 48 à 50)}

Soit $R$ une relation d'équivalence sur $S$ dont les classes $S_i$ sont deux
à deux non parallèles --- une partition non croisée. Un segment est
\emph{saturé} s'il est réunion de classes. Le complémentaire d'un segment
saturé est saturé ; la réunion et l'intersection de deux segments saturés
qui se rencontrent sans recouvrir $S$ le sont aussi. Pour une classe $S_i$,
les segments composants de $S \smallsetminus S_i$ sont saturés : c'est la
condition c) de la page 43. Leurs complémentaires, qui sont les segments
minimaux contenant $S_i$, le sont donc aussi\footnote{Les pages 48 à 50 sont
à l'encre brune, d'une main rapide, dont les formules portent et la prose
souvent pas. On a gardé ce qui s'y démontre.}.

\textbf{Proposition (page 48).} \emph{Les segments saturés minimaux sont les
classes de $R$ qui sont des segments.} Il les appelle \emph{anglets}.

Une classe qui est un segment est un segment saturé, évidemment minimal.
Réciproquement, soit $A$ un segment saturé minimal, qu'on regarde comme un
intervalle ordonné par $u$, et soit $S_k$ la classe qui contient la première
extrémité de $A$. Les éléments de $A$ qui suivent le dernier point de $S_k$
forment, s'il y en a, un segment saturé strictement contenu dans $A$ : une
classe qui aurait des points de part et d'autre de ce dernier point
croiserait $S_k$. Par minimalité, $S_k$ contient donc les deux extrémités
de $A$ ; et un segment composant de $S \smallsetminus S_k$ contenu dans $A$
serait encore un segment saturé strictement plus petit. Donc $S_k =
A$\footnote{La page raisonne sur deux classes $S_i$, $S_j$ contenues dans $A$
et sur « le segment composant » de l'une qui contient l'autre, en affirmant
qu'il est contenu dans $A$. Ce n'est pas toujours le cas : si $S_i$ et $S_j$
sont côte à côte dans $A$, le composant de $S \smallsetminus S_j$ qui contient
$S_i$ déborde de $A$. La démonstration donnée ici est de nous.}.

\textbf{Corollaire (page 49).} \emph{Tout segment saturé contient un anglet.
Si $R$ n'est pas la relation grossière, il y a au moins deux anglets.} Il y
a en effet alors deux segments saturés complémentaires : un composant de $S
\smallsetminus S_i$, pour une classe $S_i \neq S$, et son complémentaire.

C'est, pour les partitions non croisées, l'analogue du fait qu'un arbre à au
moins deux sommets a au moins deux feuilles ; le mot « anglet » dit bien
qu'une telle classe occupe un coin du polygone\footnote{L'analogie est de
nous. « Anglet » est la lecture retenue pour un mot écrit six fois.}.

\textbf{L'effeuillage (pages 49 et 50).} Soit $T_1$ la réunion des anglets
et $S_1 = S \smallsetminus T_1$. Si $S_1 \neq \emptyset$, il a au moins deux
segments composants : sinon $S_1$, complémentaire du saturé $T_1$, serait un
segment saturé, et contiendrait un anglet\footnote{La page écrit « sinon,
$R$, serait un segment contenant deux anglets, absurde ». On lit : « $S_1$
serait un segment saturé, contenant donc un anglet ». Sur cette page, $R_1$
désigne tantôt la relation induite sur $S_1$, tantôt $S_1$ lui-même ; on
écrit $S_1$ et $R_1 = R|_{S_1}$.}. Aucune classe contenue dans $S_1$ n'est
contenue dans un seul segment composant de $S_1$ : sinon, n'étant pas un
segment, elle laisserait entre deux de ses points un segment saturé contenu
dans $S_1$, qui contiendrait un anglet. Alors :
\begin{itemize}
  \item si $\operatorname{card} S_1 = 2$, $S_1$ est une classe, et $R$ est
  formée des anglets et de $S_1$ ;
  \item si $\operatorname{card} S_1 \geqslant 3$, $S_1$ muni de l'ordre
  circulaire induit est un polygone, et $R_1 = R|_{S_1}$ y est une relation
  non parallèle ; si elle est grossière, $R$ est formée des anglets et de
  $S_1$ ; sinon on recommence avec les anglets de $R_1$, qui sont au moins
  deux.
\end{itemize}
La page 50 esquisse la tour qui en résulte, une suite
$\Sigma_1, \Sigma_2, \ldots$ d'ensembles de segments saturés deux à deux
disjoints, de réunions $T_1 \subset T_2 \subset \cdots$, chaque élément de
$\Sigma_i$ contenant au moins un segment composant de $T_{i-1}$ et ayant ses
extrémités hors de $T_{i-1}$. Les conditions qu'elle écrit ne sont pas
menées à terme\footnote{Plusieurs conditions sont écrites puis biffées
($\operatorname{card} \Sigma_1 \geqslant 2$, $\operatorname{card} S_1
\geqslant 2$), et la dernière, $\operatorname{card} \pi_0(S_2) \neq 1$,
termine la page ; on n'essaie pas de reconstituer la définition visée.}.
C'est une description récursive des partitions non croisées par retrait
successif des classes d'angle, comme on élague un arbre par ses feuilles.

\pagerange{51}{52}
\subsection*{Le quasi-segment (pages 51 et 52)}

Les pages 51 et 52, à l'encre noire, posent les notions des pages 41 à 47
pour un \emph{quasi-segment} $I$ : un ensemble fini totalement ordonné (à
renversement près) d'au moins deux éléments, comme un polygone qu'on aurait
ouvert le long d'une arête\footnote{« Quasi-segment », « fini » et
« sommets » sont des lectures incertaines ; la description comme polygone
ouvert est de nous. Les notions qui suivent n'ont de sens que si $I$ est
ordonné de la sorte.}. Pour $\varepsilon \in \mathfrak{P}_2(I)$, soit
$I_\varepsilon$ le segment (l'intervalle) qu'il détermine. Deux paires
$\varepsilon$, $\varepsilon'$ sont \emph{non parallèles} si elles sont
disjointes et si $I_\varepsilon$ et $I_{\varepsilon'}$ sont disjoints ou
emboîtés ; deux parties $A$, $A'$ le sont si elles sont disjointes et que
toute paire de l'une est non parallèle à toute paire de l'autre. Cela revient
à dire que $A'$ est contenu dans l'un des composants de $I \smallsetminus A$,
ou dans la réunion des deux composants extrêmes (ceux qui contiennent les
extrémités de $I$), quand ils sont deux.

\textbf{Partitions non parallèles de $I$.} Soit $\Sigma$ un ensemble de
segments de $I$ tel que :
\begin{enumerate}
  \item[a)] deux éléments de $\Sigma$ sont disjoints ou emboîtés ;
  \item[b)] si $B \subsetneq A$ dans $\Sigma$, alors $B \cap \partial A =
  \emptyset$ ;
  \item[c)] $\partial I \subset \bigcup_{A \in \Sigma} A$, donc $\partial I
  \subset \bigcup_{A \in \Sigma} \partial A$\footnote{Le signe de b) est
  surchargé ; la lecture retenue est celle de ($\sigma_2$). Le « donc » de c)
  est juste : le plus petit élément de $\Sigma$ qui contient une extrémité de
  $I$ l'a pour extrémité.}.
\end{enumerate}
Les cœurs $A^{*}$ sont non vides et deux à deux non parallèles ; avec $I' =
\bigcup A = \bigcup A^{*}$, $Z = I \smallsetminus I'$, $Z^{*} = Z$ et
$\widehat{\Sigma}$ comme plus haut, $\widehat{\Sigma}^{*}$ est une partition
non parallèle de $I$.

\emph{NB} (page 52) : toute partition non parallèle s'obtient ainsi, et de
plusieurs façons. C'est vrai : la famille des enveloppes (plus petits
segments contenants) des classes convient, car ces enveloppes sont
disjointes ou emboîtées, une classe emboîtée dans l'enveloppe d'une autre
n'en touche pas les extrémités, et le cœur de l'enveloppe d'une classe est
cette classe. Sur $I = \{1, \ldots, 5\}$, la partition $\{1\}, \{2, 4\},
\{3\}, \{5\}$ provient aussi de $\Sigma = \{\{1\}, \{3\}, \{5\}\}$, avec $Z =
\{2, 4\}$\footnote{La vérification et l'exemple sont de nous ; la page
affirme le NB sans le démontrer.}.

Soit enfin $\Sigma_{\max}$ l'ensemble des éléments maximaux de $\Sigma$ :
deux d'entre eux sont disjoints, et $I' = \bigcup_{A \in \Sigma_{\max}} A$. La
page s'arrête là, après quatre lignes.

\pagerange{53}{54}
\section*{VII. Attachement de deux arêtes (pages 53 et 54)}

Une feuille carrée revient aux cartes. Appelons \emph{binaire} une arête
bordée par deux faces distinctes, \emph{monaire} une arête bordée des deux
côtés par la même face. On part de deux cartes connexes $M'$, $M''$ sur des
surfaces orientées compactes, de genres $g'$, $g''$, à $f'$ et $f''$ faces,
et d'une arête dans chacune. La page ne décrit pas l'opération ; ses bilans
sont ceux de l'opération suivante\footnote{La reconstitution est de nous.
Elle est soutenue par les croquis : « bi » et « bi », « mon » et « bi »,
« mon » et « mon », avec des doubles flèches entre les lobes, et, page 54, un
disque relié par un tube court à une feuille rectangulaire. La condition de
la première ligne, « $(a,\omega)$, $(a',\omega)$ et $(a,\overline{\omega})$,
$(a',\overline{\omega})$ sont des paires associées », d'une lecture
incertaine, paraît fixer quelle rive de $a$ va avec quelle rive de $a'$.} :
on coupe les deux arêtes $a = xy$ et $a' = x'y'$ en leur milieu et l'on
raccorde les moitiés en croix, de sorte qu'elles deviennent $xy'$ et $x'y$ ;
sur les surfaces, on relie les deux arêtes par un tube planté en leur
milieu. Les sommets et les arêtes s'ajoutent, $V = V' + V''$ et $E = E'
+ E''$, et seules changent les faces qui bordent $a$ et $a'$ :
\begin{itemize}
  \item \emph{deux arêtes binaires} : les quatre faces qui les bordent se
  raccordent deux à deux, « 4 faces deviennent 2 », et
  \[
  f = f' + f'' - 2, \qquad g = g' + g'' ;
  \]
  \item \emph{une arête monaire, une binaire} : « trois faces deviennent 1 »,
  et encore $f = f' + f'' - 2$, $g = g' + g''$\footnote{La page écrit ici
  $f^{*}$ là où les deux autres lignes ont $f$.} ;
  \item \emph{deux arêtes monaires}, incidentes à des faces distinctes : la
  face qui borde $a$ et celle qui borde $a'$ se raccordent en une région qui
  se recoupe en deux faces, « deux faces deviennent 2 », et
  \[
  f = f' + f'', \qquad g = g' + g'' - 1 .
  \]
\end{itemize}
Le genre se déduit de la formule d'Euler $V - E + f = 2 - 2g$ pour une carte
connexe : dans les deux premiers cas $2 - 2g = (2 - 2g') + (2 - 2g'') - 2$,
dans le troisième $2 - 2g = (2 - 2g') + (2 - 2g'')$\footnote{Dans le
troisième cas, la formule suppose que la carte obtenue est connexe, ce que
la page ne dit pas. Si $a$ et $a'$ sont deux isthmes --- arêtes dont le
retrait déconnecte, toujours monaires ---, le raccord en croix produit deux
cartes, et $g' + g'' - 1$ peut valoir $-1$ ; il faut qu'au moins l'une des
deux arêtes monaires ne soit pas un isthme, ce qui demande un genre au moins
égal à $1$. La condition est de nous.}. Dans le troisième cas, la région
formée par le tube n'est plus un disque, et la carte, redessinée sur la
surface qu'elle définit, en perd une anse.

C'est, sur les cartes combinatoires, le pendant du recollement que le
dossier 67 étudie pour les surfaces conformes (pages 2 et 3, puis 53 à
60)\footnote{Le rapprochement est de nous ; aucune des deux séries ne cite
l'autre.}. La page 54 ne porte que des croquis au crayon de la même
opération en volume.

\end{document}
